\section{高级算法专题章正文案例推导补强}

这一节补齐本章正文出现过的 LeetCode 案例推导。阅读顺序统一按“题目说明、样例入口、暴力基线、暴力过程、重复工作、优化条件、相邻状态、状态复用、指针和字段、代码落点、正确性、边界实验、复杂度变化、迁移追问”展开；目的不是新增题量，而是把正文中的案例都讲成可复述、可证明、可迁移的解题链。

\subsection{LeetCode 731 我的日程安排表 II}

\textbf{题目说明。} LeetCode 731 我的日程安排表 II 的题面入口要先还原成具体任务：允许双重预订但不允许三重预订，插入会改变区间重叠层数。

\textbf{样例入口。} 拿 book(10, 20)、book(15, 25) 后，重叠段 [15, 20) 已经是双重预订；再 book(17, 22) 会碰到双重段，必须失败。

\textbf{暴力基线。} 慢版本对新区间和所有旧区间两两比较，手工判断相交、覆盖或冲突。它能验证端点语义，但排序后很多远离当前边界的区间无需再看。放到本题就是：先按“允许双重预订但不允许三重预订，插入会改变区间重叠层数”完整试慢版本；样例里已经能看到“规律是保存所有单次预订和已形成的双重预订，新区间不能与任何双重区间相交”，所以优化前必须先能说清慢版本枚举了哪些候选。

\textbf{暴力过程。} 拿 book(10, 20)、book(15, 25) 后，重叠段 [15, 20) 已经是双重预订；再 book(17, 22) 会碰到双重段，必须失败。

\textbf{重复工作。} 题面模型：允许双重预订但不允许三重预订，插入会改变区间重叠层数。暴力版的慢点要落到本题：慢版本会围绕“允许双重预订但不允许三重预订，插入会改变区间重叠层数”做两两区间比较。样例规律是“规律是保存所有单次预订和已形成的双重预订，新区间不能与任何双重区间相交”，说明排序后只有贴近当前端点、结果尾部或资源堆顶的区间还会影响答案。后面的优化只消掉这类重复工作，不能改变答案集合。

\textbf{优化条件。} 本题的关键观察是：保存已有预订和已经双重预订的区间；新预订若碰到双重区间则失败，否则记录新产生的重叠。这句话必须能翻译成可维护状态；如果题目条件改变后观察不再成立，就不能继续套原来的写法。

\textbf{相邻状态。} 规律是保存所有单次预订和已形成的双重预订，新区间不能与任何双重区间相交。把这个特例继续拆开看：排序以后，真正还能影响当前区间的候选必须贴近当前端点、结果尾部或资源堆顶。某个区间一旦被覆盖、合并或弹出，就要说明它为什么不会和后续区间重新产生更优关系。

\textbf{状态复用。} 把两两比较压成排序后的相邻关系、当前边界或动态有序结构；远离当前边界的区间要被安全归档。本题落点是：规律是保存所有单次预订和已形成的双重预订，新区间不能与任何双重区间相交。手算时只改被新输入影响的那一小部分，而不是回头重算完整候选。

\textbf{指针和字段。} 状态字段要能说清：排序后的区间、当前结果尾、扫描右端或资源堆、端点比较规则；动态版本还要保存前驱后继。手算时按这个协议跟踪：拿 book(10, 20)、book(15, 25) 后，重叠段 [15, 20) 已经是双重预订；再 book(17, 22) 会碰到双重段，必须失败。然后按协议复查：手算时先排序，再逐个区间写当前结果尾部、当前右端或资源堆。每个区间离开视野时都要说明它不会再影响后续。

\textbf{代码落点。} 排序比较器满足严格弱序；端点相等的处理和题意一致；扫描时只和当前结果尾或当前资源集合交互。关键代码行必须能逐句翻译成“旧状态怎样变成新状态”，否则就只是模板。

\textbf{正确性。} 证明从排序后的邻接关系开始：已经合并、选择或丢弃的区间不会被更后面的区间重新影响。

\textbf{边界实验。} 先用端点相接、多个局部重叠、完全覆盖、回滚失败插入做反例检查。实验时覆盖端点相接、多个局部重叠、完全覆盖、回滚失败插入，画出排序后的区间序列和当前结果尾部。动态版本要专门测前驱、后继和端点相接。

\textbf{复杂度变化。} 暴力两两比较区间常见为平方级；排序后扫描通常由排序成本主导，动态有序结构则把每次前驱后继查询控制在对数级。

\textbf{迁移追问。} 如果题目改成“扫描线差分也能做，但要处理失败时恢复状态”，先判断原状态是否仍然覆盖新答案集合，再决定是微调边界、改 value 语义，还是换算法。

\subsection{LeetCode 1483 树节点的第 K 个祖先}

\textbf{题目说明。} LeetCode 1483 树节点的第 K 个祖先 的题面入口要先还原成具体任务：第 k 个祖先可以按 k 的二进制拆成若干 \ensuremath{2^{j}} 级跳跃。

\textbf{样例入口。} 拿一棵树中查询 node=5 的第 6 个祖先手算，6 的二进制是 110，可以先跳 2 级再跳 4 级，而不是一步步向上走 6 次。

\textbf{暴力基线。} 慢版本在每个节点重新扫描子树或路径，先求出局部答案。重复扫描暴露了真正状态：每个节点只需要从孩子或父亲那里拿到少量信息。放到本题就是：先按“第 k 个祖先可以按 k 的二进制拆成若干 \ensuremath{2^{j}} 级跳跃”完整试慢版本；样例里已经能看到“规律是预处理 up[node][j] 表示 \ensuremath{2^{j}} 祖先，查询时按 k 的二进制位跳跃”，所以优化前必须先能说清慢版本枚举了哪些候选。

\textbf{暴力过程。} 拿一棵树中查询 node=5 的第 6 个祖先手算，6 的二进制是 110，可以先跳 2 级再跳 4 级，而不是一步步向上走 6 次。

\textbf{重复工作。} 题面模型：第 k 个祖先可以按 k 的二进制拆成若干 \ensuremath{2^{j}} 级跳跃。暴力版的慢点要落到本题：慢版本会在树上重复确认“第 k 个祖先可以按 k 的二进制拆成若干 \ensuremath{2^{j}} 级跳跃”。样例规律是“规律是预处理 up[node][j] 表示 \ensuremath{2^{j}} 祖先，查询时按 k 的二进制位跳跃”，说明父子之间真正需要传递的是高度、上下界、命中状态、层号或两类选择值，而不是反复扫描整棵子树。后面的优化只消掉这类重复工作，不能改变答案集合。

\textbf{优化条件。} 本题的关键观察是：预处理 up[node][j] 表示 node 的 \ensuremath{2^{j}} 祖先，查询时按 k 的二进制位向上跳。这句话必须能翻译成可维护状态；如果题目条件改变后观察不再成立，就不能继续套原来的写法。

\textbf{相邻状态。} 规律是预处理 up[node][j] 表示 \ensuremath{2^{j}} 祖先，查询时按 k 的二进制位跳跃。把这个特例继续拆开看：节点向父亲返回的量和全局答案通常不是同一个量。先写叶子或空节点的返回值，再写父节点如何合并左右孩子，递归式才有边界基础。

\textbf{状态复用。} 把对子树的重复扫描压成返回值或队列状态；每个节点只合并孩子给出的稳定信息。本题落点是：规律是预处理 up[node][j] 表示 \ensuremath{2^{j}} 祖先，查询时按 k 的二进制位跳跃。手算时只改被新输入影响的那一小部分，而不是回头重算完整候选。

\textbf{指针和字段。} 状态字段要能说清：节点指针、传入约束或返回值、当前答案、层号或路径状态；空节点的返回值必须和状态定义匹配。手算时按这个协议跟踪：拿一棵树中查询 node=5 的第 6 个祖先手算，6 的二进制是 110，可以先跳 2 级再跳 4 级，而不是一步步向上走 6 次。然后按协议复查：手算时从叶子开始写返回值，或者从根开始写传入约束。不要只写访问顺序，要写每条边上传递的信息。

\textbf{代码落点。} 递归版本先处理空节点；迭代版本的栈元素要带完整状态；全负值、重复值和极端深树要单独测。关键代码行必须能逐句翻译成“旧状态怎样变成新状态”，否则就只是模板。

\textbf{正确性。} 证明从子树归纳或层序不变量开始：孩子状态正确时父节点合并正确；若用队列，当前轮队列必须正好保存当前层节点。

\textbf{边界实验。} 先用根节点祖先不存在、k 为零、树很深、查询很多、表大小做反例检查。实验时覆盖根节点祖先不存在、k 为零、树很深、查询很多、表大小，尤其关注空节点、单节点、链式树和全负值。递归版本和显式栈版本可以互相对拍。

\textbf{复杂度变化。} 暴力在每个节点重复扫描子树或反复寻找层边界可能退化到平方级；后序汇总、显式栈或层序遍历让每个节点处理常数次，通常是线性时间。

\textbf{迁移追问。} 如果题目改成“LCA 也使用倍增表，只是同时提升两个节点”，先判断原状态是否仍然覆盖新答案集合，再决定是微调边界、改 value 语义，还是换算法。

\subsection{LeetCode 212 单词搜索 II}

\textbf{题目说明。} LeetCode 212 单词搜索 II 的题面入口要先还原成具体任务：多个单词在同一棋盘上搜索，公共前缀可以共享。

\textbf{样例入口。} 拿 board 上有 oath、pea、eat 手算，单词共享前缀时不应对每个词重复全图搜索；从 o 出发沿 Trie 前缀 o-a-t-h 找到 oath。

\textbf{暴力基线。} 慢版本就是完整搜索树：每层列出选择列表，路径到达终止条件时检查答案。它先保证覆盖所有可能，再把明显无效或重复的分支剪掉。放到本题就是：先按“多个单词在同一棋盘上搜索，公共前缀可以共享”完整试慢版本；样例里已经能看到“规律是 Trie 合并所有单词前缀，DFS 走棋盘同时走 Trie，命中单词后记录并可去重”，所以优化前必须先能说清慢版本枚举了哪些候选。

\textbf{暴力过程。} 拿 board 上有 oath、pea、eat 手算，单词共享前缀时不应对每个词重复全图搜索；从 o 出发沿 Trie 前缀 o-a-t-h 找到 oath。

\textbf{重复工作。} 题面模型：多个单词在同一棋盘上搜索，公共前缀可以共享。暴力版的慢点要落到本题：慢版本会围绕“多个单词在同一棋盘上搜索，公共前缀可以共享”展开完整搜索树。样例规律是“规律是 Trie 合并所有单词前缀，DFS 走棋盘同时走 Trie，命中单词后记录并可去重”，说明一层递归必须对应一个明确决策，剪枝只能删除整棵不可能成功或重复的子树。后面的优化只消掉这类重复工作，不能改变答案集合。

\textbf{优化条件。} 本题的关键观察是：先建 Trie，再从每个格子回溯；路径到达单词结束节点就收集答案。这句话必须能翻译成可维护状态；如果题目条件改变后观察不再成立，就不能继续套原来的写法。

\textbf{相邻状态。} 规律是 Trie 合并所有单词前缀，DFS 走棋盘同时走 Trie，命中单词后记录并可去重。把这个特例继续拆开看：一层递归对应一个明确决策，路径保存已经做出的选择，选择列表保存仍可能扩展的分支。剪枝前必须能说明被剪掉的整棵子树没有合法答案，而不是只因为它看起来不优。

\textbf{状态复用。} 把完整搜索树加上合法性检查、剩余资源剪枝和同层去重；路径状态进入和退出要成对恢复。本题落点是：规律是 Trie 合并所有单词前缀，DFS 走棋盘同时走 Trie，命中单词后记录并可去重。手算时只改被新输入影响的那一小部分，而不是回头重算完整候选。

\textbf{指针和字段。} 状态字段要能说清：路径、起点或使用标记、剩余目标、答案集合、剪枝条件；进入递归和退出递归必须成对恢复。手算时按这个协议跟踪：拿 board 上有 oath、pea、eat 手算，单词共享前缀时不应对每个词重复全图搜索；从 o 出发沿 Trie 前缀 o-a-t-h 找到 oath。然后按协议复查：手算时给每层标出选择列表和路径。剪枝发生前要指出被剪分支为什么已经不可能得到合法答案。

\textbf{代码落点。} 剪枝条件写在扩展前；同层去重不要误写成全局去重；保存答案时复制当前路径。关键代码行必须能逐句翻译成“旧状态怎样变成新状态”，否则就只是模板。

\textbf{正确性。} 证明从搜索树覆盖开始：每个合法答案有且只有一条路径，剪枝只删掉不可能成功的分支。

\textbf{边界实验。} 先用重复单词、同一格不能复用、前缀是完整单词、剪枝删除空子树做反例检查。实验时覆盖重复单词、同一格不能复用、前缀是完整单词、剪枝删除空子树，统计搜索节点数量和剪枝次数。重复元素题要打印同一层跳过哪些值，确认没有误删合法路径。

\textbf{复杂度变化。} 暴力搜索树的规模由输出或选择空间决定；剪枝不改变最坏指数级本质，但能删除不可能成功或重复的分支，复杂度要说明输出规模。

\textbf{迁移追问。} 如果题目改成“这是单词搜索从单目标扩展到多目标后的结构升级”，先判断原状态是否仍然覆盖新答案集合，再决定是微调边界、改 value 语义，还是换算法。

\subsection{LeetCode 1192 查找集群内的关键连接}

\textbf{题目说明。} LeetCode 1192 查找集群内的关键连接 的题面入口要先还原成具体任务：无向图中的桥是删除后会断开连通性的边。

\textbf{样例入口。} 拿 n=4、connections=[[0, 1], [1, 2], [2, 0], [1, 3]] 手算，0-1-2 形成环，删除其中任一条仍可互通；节点 3 只通过边 1-3 接入，删除它会断开。

\textbf{暴力基线。} 慢版本先按题意画出完整状态图，用普通搜索跑通可达性。这个过程用于确认无向图中的桥是删除后会断开连通性的边的节点和边，之后再讨论访问标记、层次或代价顺序。放到本题就是：先按“无向图中的桥是删除后会断开连通性的边”完整试慢版本；样例里已经能看到“规律是 DFS 中 low[child] 若大于 disc[parent]，说明 child 子树没有备用返祖通道，这条边就是桥”，所以优化前必须先能说清慢版本枚举了哪些候选。

\textbf{暴力过程。} 拿 n=4、connections=[[0, 1], [1, 2], [2, 0], [1, 3]] 手算，0-1-2 形成环，删除其中任一条仍可互通；节点 3 只通过边 1-3 接入，删除它会断开。

\textbf{重复工作。} 题面模型：无向图中的桥是删除后会断开连通性的边。暴力版的慢点要落到本题：慢版本会围绕“无向图中的桥是删除后会断开连通性的边”反复展开同一批状态。样例规律是“规律是 DFS 中 low[child] 若大于 disc[parent]，说明 child 子树没有备用返祖通道，这条边就是桥”，说明必须把节点、边、访问标记和资源维度一次建清；同一状态不应重复入队，不同资源状态也不能误合并。后面的优化只消掉这类重复工作，不能改变答案集合。

\textbf{优化条件。} 本题的关键观察是：DFS 记录 disc 和 low，若 low[child] 大于 disc[parent]，子树没有返祖通道，父子边是桥。这句话必须能翻译成可维护状态；如果题目条件改变后观察不再成立，就不能继续套原来的写法。

\textbf{相邻状态。} 规律是 DFS 中 low[child] 若大于 disc[parent]，说明 child 子树没有备用返祖通道，这条边就是桥。把这个特例继续拆开看：状态节点不一定等于原始位置，还可能包含钥匙、剩余资源、时间或访问集合。visited 只能合并未来能力完全相同的状态，否则会把合法路径剪掉。

\textbf{状态复用。} 把重复搜索压成访问标记、距离数组或拓扑状态；邻居生成也要从全量扫描变成结构化枚举。本题落点是：规律是 DFS 中 low[child] 若大于 disc[parent]，说明 child 子树没有备用返祖通道，这条边就是桥。手算时只改被新输入影响的那一小部分，而不是回头重算完整候选。

\textbf{指针和字段。} 状态字段要能说清：状态编号、邻接生成方式、访问标记、距离或层数、队列内容；若状态含资源，visited 不能只按位置记录。手算时按这个协议跟踪：拿 n=4、connections=[[0, 1], [1, 2], [2, 0], [1, 3]] 手算，0-1-2 形成环，删除其中任一条仍可互通；节点 3 只通过边 1-3 接入，删除它会断开。然后按协议复查：手算时画队列或栈，状态必须包含题目需要的所有资源。入队时标记还是出队时标记，要和去重语义一致。

\textbf{代码落点。} 入队时标记还是出队时标记必须一致；方向数组集中定义；多源 BFS 要一次性把所有源点放入初始队列。关键代码行必须能逐句翻译成“旧状态怎样变成新状态”，否则就只是模板。

\textbf{正确性。} 证明从可达性开始：搜索覆盖所有合法边能到达的状态，访问标记只删除重复状态而不删除不同未来能力。

\textbf{边界实验。} 先用重边、根节点、图保证连通、无向边需要编号跳过反向边、递归栈做反例检查。实验时覆盖重边、根节点、图保证连通、无向边需要编号跳过反向边、递归栈，并打印入队时的完整状态。若同一位置但资源不同的状态被合并，很多图题会出现隐蔽错误。

\textbf{复杂度变化。} 暴力重复从多个起点搜索会反复遍历图；正确建模和访问标记后，BFS 或 DFS 通常按节点数加边数计算，状态图题则按完整状态数计算。

\textbf{迁移追问。} 如果题目改成“若要找割点，判断条件和根节点处理不同”，先判断原状态是否仍然覆盖新答案集合，再决定是微调边界、改 value 语义，还是换算法。

\subsection{LeetCode 307 区域和检索 - 数组可修改}

\textbf{题目说明。} LeetCode 307 区域和检索 - 数组可修改 的题面入口要先还原成具体任务：单点更新会破坏普通前缀和后面所有位置。

\textbf{样例入口。} 拿 nums=[1, 3, 5] 手算，sumRange(0, 2)=9；update(1, 2) 后数组变成 [1, 2, 5]，后续区间和都受影响。普通前缀和要改很多位置，树状数组只把差值 -1 加到若干 lowbit 管辖节点。

\textbf{暴力基线。} 普通前缀和能 O(1) 查询，但 update(index, val) 后 index 之后所有前缀都要修改。

\textbf{暴力过程。} 以 nums=[1, 3, 5] 为例，update(1, 2) 让数组变成 [1, 2, 5]。若维护 prefix，prefix[2]、prefix[3] 都要改。

\textbf{重复工作。} 题面模型：单点更新会破坏普通前缀和后面所有位置。暴力版的慢点要落到本题：浪费在一次单点更新影响大量前缀。树状数组把前缀和拆成若干 lowbit 管辖段，更新只触碰包含该点的少数段。后面的优化只消掉这类重复工作，不能改变答案集合。

\textbf{优化条件。} 本题的关键观察是：树状数组把前缀和拆成若干 lowbit 管辖段，更新和查询都只访问对数个节点。这句话必须能翻译成可维护状态；如果题目条件改变后观察不再成立，就不能继续套原来的写法。

\textbf{相邻状态。} 规律是树状数组用对数个节点维护前缀和，更新和查询都沿 lowbit 跳转。把这个特例继续拆开看：先标出已经确认的前缀或后缀，再标出未知区；能被丢掉的候选，必须是以后再也不会改变已确认区域语义的候选。这样才算从例子推出数组不变量，而不是只记一次操作。

\textbf{状态复用。} BIT 使用 1 基下标。update 计算 delta，然后 for i+=lowbit(i) 更新管辖段；query 前缀和时 for i-=lowbit(i) 累加。手算时只改被新输入影响的那一小部分，而不是回头重算完整候选。

\textbf{指针和字段。} 状态字段要能说清：tree[i] 保存一段长度 lowbit(i) 的区间和，nums 保存当前真实值以计算 delta。手算时按这个协议跟踪：update(1, 2) 的 1 基位置是 2，delta=-1，更新 tree[2]、tree[4] 等管辖节点。

\textbf{代码落点。} 关键代码是 index += index \& -index 和 index -= index \& -index；sumRange(l, r)=prefix(r+1)-prefix(l)。关键代码行必须能逐句翻译成“旧状态怎样变成新状态”，否则就只是模板。

\textbf{正确性。} 任意前缀可以被若干 lowbit 段不重叠拆分；单点更新只需更新所有包含该点的拆分段。

\textbf{边界实验。} 先用下标一基转换、负数更新、重复更新、查询单点、差值更新做反例检查。实验时把下标一基转换、负数更新、重复更新、查询单点、差值更新列成表，再打印关键下标和有效区间。若某个元素被写入后又被未来步骤破坏，说明区间不变量没有成立。

\textbf{复杂度变化。} 普通前缀 update O(n)，树状数组 update/query 均 O(log n)，空间 O(n)。

\textbf{迁移追问。} 如果题目改成“若需要区间更新或更复杂信息，线段树更通用”，先判断原状态是否仍然覆盖新答案集合，再决定是微调边界、改 value 语义，还是换算法。

\subsection{LeetCode 315 计算右侧小于当前元素的个数}

\textbf{题目说明。} LeetCode 315 计算右侧小于当前元素的个数 的题面入口要先还原成具体任务：从右向左扫描时，已处理元素正好是当前位置右侧元素集合。

\textbf{样例入口。} 拿 nums=[5, 2, 6, 1] 从右往左手算，先加入 1；处理 6 时右侧小于 6 的已有元素有 1 个；处理 2 时右侧小于 2 的也只有 1；处理 5 时右侧小于 5 的是 2 和 1。

\textbf{暴力基线。} 慢版本按顺序扫描下标、逐个试答案值，或直接检查候选直到遇到目标边界。它先保证候选空间覆盖完整，但每次只排除一个候选，浪费了题目给出的顺序、比较或可行性边界。放到本题就是：先按“从右向左扫描时，已处理元素正好是当前位置右侧元素集合”完整试慢版本；样例里已经能看到“规律是右侧元素集合动态增长，离散化后用树状数组查询当前排名之前的频次”，所以优化前必须先能说清慢版本枚举了哪些候选。

\textbf{暴力过程。} 拿 nums=[5, 2, 6, 1] 从右往左手算，先加入 1；处理 6 时右侧小于 6 的已有元素有 1 个；处理 2 时右侧小于 2 的也只有 1；处理 5 时右侧小于 5 的是 2 和 1。

\textbf{重复工作。} 题面模型：从右向左扫描时，已处理元素正好是当前位置右侧元素集合。暴力版的慢点要落到本题：慢版本会按顺序试“从右向左扫描时，已处理元素正好是当前位置右侧元素集合”里的候选；而样例规律已经指出“规律是右侧元素集合动态增长，离散化后用树状数组查询当前排名之前的频次”。一次比较或一次可行性检查能丢掉一整段候选，继续线性试就是浪费。后面的优化只消掉这类重复工作，不能改变答案集合。

\textbf{优化条件。} 本题的关键观察是：先离散化数值，再用树状数组维护右侧元素排名频次，查询当前排名之前的频次和。这句话必须能翻译成可维护状态；如果题目条件改变后观察不再成立，就不能继续套原来的写法。

\textbf{相邻状态。} 规律是右侧元素集合动态增长，离散化后用树状数组查询当前排名之前的频次。把这个特例继续拆开看：每次比较 mid 之后，必须能指出哪一侧整体不可能包含目标，或哪一侧整体已经满足可行性。二分的核心不是 mid 公式，而是被丢弃区间确实不含答案。

\textbf{状态复用。} 把逐个试候选改成维护答案所在区间；边界谓词、可行性检查或局部比较给出分支方向，每个分支都必须能排除一侧候选。本题落点是：规律是右侧元素集合动态增长，离散化后用树状数组查询当前排名之前的频次。手算时只改被新输入影响的那一小部分，而不是回头重算完整候选。

\textbf{指针和字段。} 状态字段要能说清：left、right、mid、mid 处比较值或检查返回值、答案区间含义、返回值语义；每一轮更新都要保留目标位置或第一个可行值。手算时按这个协议跟踪：拿 nums=[5, 2, 6, 1] 从右往左手算，先加入 1；处理 6 时右侧小于 6 的已有元素有 1 个；处理 2 时右侧小于 2 的也只有 1；处理 5 时右侧小于 5 的是 2 和 1。然后按协议复查：手算时列出 left、mid、right，以及 mid 处比较结果、可行性检查结果或局部趋势。每一步都要写出被丢弃的一侧为什么不含答案。

\textbf{代码落点。} 检查函数或比较分支单独写清，mid 不溢出，循环退出条件和返回值配套；每个分支都要解释丢弃区间。关键代码行必须能逐句翻译成“旧状态怎样变成新状态”，否则就只是模板。

\textbf{正确性。} 证明从排除一侧开始：答案空间题证明谓词单调，局部比较题证明保留下来的区间仍含目标；不能只靠经验猜测左右边界。

\textbf{边界实验。} 先用重复值、负数、值域很大、严格小于不能包含相等值、结果顺序要反向写回做反例检查。实验时覆盖重复值、负数、值域很大、严格小于不能包含相等值、结果顺序要反向写回，逐轮打印 left、mid、right、比较值或检查结果。只要有一次被排除区间仍含答案，二分不变量就错了。

\textbf{复杂度变化。} 暴力逐个试候选是线性或和值域成正比；二分每次排除一半候选，复杂度变成对数级，答案二分还要乘上单次检查成本。

\textbf{迁移追问。} 如果题目改成“若改成右侧大于当前元素，只需查询排名之后的频次”，先判断原状态是否仍然覆盖新答案集合，再决定是微调边界、改 value 语义，还是换算法。

\subsection{LeetCode 493 翻转对}

\textbf{题目说明。} LeetCode 493 翻转对 的题面入口要先还原成具体任务：每个 i 需要统计右侧满足 nums[i] 大于 2*nums[j] 的元素。

\textbf{样例入口。} 拿 nums=[1, 3, 2, 3, 1] 手算，翻转对有 (3, 1)、(3, 1) 两个。暴力每个 i 扫右侧是平方级；归并排序中左右半区已有序，可以用指针统计左侧每个数能压过右侧多少个 2*nums[j]。

\textbf{暴力基线。} 暴力枚举所有 i<j，检查 nums[i] > 2*nums[j]，计数满足条件的对。

\textbf{暴力过程。} 以 [1, 3, 2, 3, 1] 为例，暴力会让每个左侧元素扫描右侧所有元素，找到两个翻转对。

\textbf{重复工作。} 题面模型：每个 i 需要统计右侧满足 nums[i] 大于 2*nums[j] 的元素。暴力版的慢点要落到本题：浪费在右侧反复扫描。归并排序中左右半区已经有序，可以用指针一次统计每个左元素能匹配多少右元素。后面的优化只消掉这类重复工作，不能改变答案集合。

\textbf{优化条件。} 本题的关键观察是：归并排序在合并前用两个有序半区双指针计数，再正常归并保持有序。这句话必须能翻译成可维护状态；如果题目条件改变后观察不再成立，就不能继续套原来的写法。

\textbf{相邻状态。} 规律是计数发生在归并前，随后归并保持区间有序。把这个特例继续拆开看：先标出已经确认的前缀或后缀，再标出未知区；能被丢掉的候选，必须是以后再也不会改变已确认区域语义的候选。这样才算从例子推出数组不变量，而不是只记一次操作。

\textbf{状态复用。} 分治排序。合并前，left 半区和 right 半区都有序；对每个 i，推进 j 直到 nums[i] <= 2*nums[j]，新增 j-mid-1 个对。手算时只改被新输入影响的那一小部分，而不是回头重算完整候选。

\textbf{指针和字段。} 状态字段要能说清：i 是左半区位置，j 是右半区第一个不满足的位置，merge 负责保持当前区间有序。手算时按这个协议跟踪：[1, 3, 2, 3, 1] 分治后，在合并有序半区时统计 3 与右侧 1 的关系，严格大于 2*1 成立。

\textbf{代码落点。} 关键是比较 1LL*nums[i] > 2LL*nums[j]，避免乘以 2 溢出；计数完成后仍要正常归并排序。关键代码行必须能逐句翻译成“旧状态怎样变成新状态”，否则就只是模板。

\textbf{正确性。} 右半区有序时，对固定 i，满足条件的右侧下标形成前缀；j 单调前进不会漏计也不会重复计。

\textbf{边界实验。} 先用负数、乘以二溢出、重复值、严格大于、答案数量很大做反例检查。实验时把负数、乘以二溢出、重复值、严格大于、答案数量很大列成表，再打印关键下标和有效区间。若某个元素被写入后又被未来步骤破坏，说明区间不变量没有成立。

\textbf{复杂度变化。} 暴力 O(\ensuremath{n^{2}})，归并计数 O(n log n)，辅助空间 O(n)。

\textbf{迁移追问。} 如果题目改成“也可用离散化加树状数组，关键仍是动态排名计数”，先判断原状态是否仍然覆盖新答案集合，再决定是微调边界、改 value 语义，还是换算法。

\subsection{LeetCode 50 Pow(x, n)}

\textbf{题目说明。} LeetCode 50 Pow(x, n) 的题面入口要先还原成具体任务：指数 n 可以按二进制拆分，重复乘 x 的线性过程存在大量倍增结构。

\textbf{样例入口。} 拿 x=2、n=10 手算，10 的二进制是 1010，所以答案是 \ensuremath{2^{8}} 乘 \ensuremath{2^{2}}。逐次乘 10 次太慢；每轮把底数平方、指数右移，当前位为 1 时乘入答案。

\textbf{暴力基线。} 慢版本先用集合、数组或布尔表保存状态，逐步执行题目动作。确认每个布尔字段的含义后，才能把它压缩成位，而不是直接写位运算公式。放到本题就是：先按“指数 n 可以按二进制拆分，重复乘 x 的线性过程存在大量倍增结构”完整试慢版本；样例里已经能看到“规律是指数按二进制拆分，负指数先转成倒数并用更宽整数处理”，所以优化前必须先能说清慢版本枚举了哪些候选。

\textbf{暴力过程。} 拿 x=2、n=10 手算，10 的二进制是 1010，所以答案是 \ensuremath{2^{8}} 乘 \ensuremath{2^{2}}。逐次乘 10 次太慢；每轮把底数平方、指数右移，当前位为 1 时乘入答案。

\textbf{重复工作。} 题面模型：指数 n 可以按二进制拆分，重复乘 x 的线性过程存在大量倍增结构。暴力版的慢点要落到本题：慢版本会围绕“指数 n 可以按二进制拆分，重复乘 x 的线性过程存在大量倍增结构”使用集合、计数或完整状态表。样例规律是“规律是指数按二进制拆分，负指数先转成倒数并用更宽整数处理”，说明哪些信息可以被逐位抵消、压缩或复用；位运算只是编码，不是证明本身。后面的优化只消掉这类重复工作，不能改变答案集合。

\textbf{优化条件。} 本题的关键观察是：快速幂每轮根据最低位决定是否乘入答案，同时底数平方、指数右移。这句话必须能翻译成可维护状态；如果题目条件改变后观察不再成立，就不能继续套原来的写法。

\textbf{相邻状态。} 规律是指数按二进制拆分，负指数先转成倒数并用更宽整数处理。把这个特例继续拆开看：每一位或每个掩码位都要对应题目里的一个布尔事实。位运算只是压缩表示；如果两个状态在位置相同但掩码不同，未来能力不同，就不能合并。

\textbf{状态复用。} 把集合状态压成掩码；包含、加入、删除、枚举子集都变成位操作，但每一位的题目含义不能丢。本题落点是：规律是指数按二进制拆分，负指数先转成倒数并用更宽整数处理。手算时只改被新输入影响的那一小部分，而不是回头重算完整候选。

\textbf{指针和字段。} 状态字段要能说清：掩码含义、当前位、目标位集合、状态转移和答案读取；负数或高位参与时要说明整数宽度。手算时按这个协议跟踪：拿 x=2、n=10 手算，10 的二进制是 1010，所以答案是 \ensuremath{2^{8}} 乘 \ensuremath{2^{2}}。逐次乘 10 次太慢；每轮把底数平方、指数右移，当前位为 1 时乘入答案。然后按协议复查：手算时把整数写成二进制，并在每一位旁边标注含义。状态转移只允许改变题目动作允许改变的那些位。

\textbf{代码落点。} 移位表达式加括号；使用无符号或足够宽类型；枚举子集时确认空集和全集是否包含。关键代码行必须能逐句翻译成“旧状态怎样变成新状态”，否则就只是模板。

\textbf{正确性。} 证明从逐位独立或子集归纳开始：被压缩到位里的信息足以决定未来转移。

\textbf{边界实验。} 先用n 为零、n 为负、int 最小值取反溢出、x 为零、浮点误差做反例检查。实验时覆盖n 为零、n 为负、int 最小值取反溢出、x 为零、浮点误差，把数字写成二进制逐位检查。负数、最高位、空掩码和全集掩码都要单独测。

\textbf{复杂度变化。} 暴力用集合或数组保存状态可能需要更多空间和比较成本；位运算通常把每步状态更新压成常数或按位数计算，掩码 DP 仍按状态数增长。

\textbf{迁移追问。} 如果题目改成“矩阵快速幂和模幂都是同一二进制拆分思想”，先判断原状态是否仍然覆盖新答案集合，再决定是微调边界、改 value 语义，还是换算法。

\subsection{LeetCode 204 计数质数}

\textbf{题目说明。} LeetCode 204 计数质数 的题面入口要先还原成具体任务：逐个数试除会重复判断很多合数的因子。

\textbf{样例入口。} 拿 n=10 手算，2 是质数，标记 4、6、8；3 是质数，标记 9；4 已被标记跳过，剩下质数是 2、3、5、7。

\textbf{暴力基线。} 暴力逐个判断 2..n-1 是否为质数，每个数再试除 2..sqrt(x)。

\textbf{暴力过程。} 以 n=10 为例，暴力会分别判断 4、6、8 是否能被 2 整除，也会在多个合数上重复试同一个因子。

\textbf{重复工作。} 题面模型：逐个数试除会重复判断很多合数的因子。暴力版的慢点要落到本题：浪费在反复证明合数是合数。一个数一旦被某个质数标记为倍数，后续不必再试除判断。后面的优化只消掉这类重复工作，不能改变答案集合。

\textbf{优化条件。} 本题的关键观察是：埃氏筛从每个质数的平方开始标记它的倍数，未被标记的数就是质数。这句话必须能翻译成可维护状态；如果题目条件改变后观察不再成立，就不能继续套原来的写法。

\textbf{相邻状态。} 规律是埃氏筛只从质数的平方开始标记倍数，避免小倍数被重复处理。把这个特例继续拆开看：先标出已经确认的前缀或后缀，再标出未知区；能被丢掉的候选，必须是以后再也不会改变已确认区域语义的候选。这样才算从例子推出数组不变量，而不是只记一次操作。

\textbf{状态复用。} is\_prime 或 is\_composite 数组记录标记状态。从质数 i 的 i*i 开始标记 i 的倍数。手算时只改被新输入影响的那一小部分，而不是回头重算完整候选。

\textbf{指针和字段。} 状态字段要能说清：i 是当前候选质数，multiple 是被标记的合数，count 是未被标记的质数数量。手算时按这个协议跟踪：n=10 时，i=2 标记 4、6、8；i=3 从 9 开始标记；4 已被标记跳过，质数为 2、3、5、7。

\textbf{代码落点。} 关键是 for (long long j=1LL*i*i; j<n; j+=i) 标记，边界是不包含 n。关键代码行必须能逐句翻译成“旧状态怎样变成新状态”，否则就只是模板。

\textbf{正确性。} 小于 i*i 的 i 的倍数已经有更小因子标记过；未被任何小因子标记的数只能是质数。

\textbf{边界实验。} 先用n 小于二、从 i*i 开始防溢出、重复标记、边界是不包含 n做反例检查。实验时把n 小于二、从 i*i 开始防溢出、重复标记、边界是不包含 n列成表，再打印关键下标和有效区间。若某个元素被写入后又被未来步骤破坏，说明区间不变量没有成立。

\textbf{复杂度变化。} 试除法约 O(n sqrt n)，埃氏筛 O(n log log n) 时间、O(n) 空间。

\textbf{迁移追问。} 如果题目改成“若要多次查询前缀质数个数，可以在筛后再做前缀计数”，先判断原状态是否仍然覆盖新答案集合，再决定是微调边界、改 value 语义，还是换算法。

\subsection{LeetCode 372 超级次方}

\textbf{题目说明。} LeetCode 372 超级次方 的题面入口要先还原成具体任务：大指数以十进制数组给出，不能先转成普通整数。

\textbf{样例入口。} 拿 a=2、b=[1, 0] 手算，目标是 \ensuremath{2^{10}} mod 1337，不能把很长指数转成普通整数。扫描十进制指数时，新答案等于旧答案的十次方乘 a 的当前位次方。

\textbf{暴力基线。} 慢版本先用集合、数组或布尔表保存状态，逐步执行题目动作。确认每个布尔字段的含义后，才能把它压缩成位，而不是直接写位运算公式。放到本题就是：先按“大指数以十进制数组给出，不能先转成普通整数”完整试慢版本；样例里已经能看到“规律是把大指数按十进制位折叠，每步都用快速幂并取模”，所以优化前必须先能说清慢版本枚举了哪些候选。

\textbf{暴力过程。} 拿 a=2、b=[1, 0] 手算，目标是 \ensuremath{2^{10}} mod 1337，不能把很长指数转成普通整数。扫描十进制指数时，新答案等于旧答案的十次方乘 a 的当前位次方。

\textbf{重复工作。} 题面模型：大指数以十进制数组给出，不能先转成普通整数。暴力版的慢点要落到本题：慢版本会围绕“大指数以十进制数组给出，不能先转成普通整数”使用集合、计数或完整状态表。样例规律是“规律是把大指数按十进制位折叠，每步都用快速幂并取模”，说明哪些信息可以被逐位抵消、压缩或复用；位运算只是编码，不是证明本身。后面的优化只消掉这类重复工作，不能改变答案集合。

\textbf{优化条件。} 本题的关键观察是：利用 \ensuremath{(a^10)^{rest}} 和最后一位指数，把指数数组逐位折叠，并在每步取模。这句话必须能翻译成可维护状态；如果题目条件改变后观察不再成立，就不能继续套原来的写法。

\textbf{相邻状态。} 规律是把大指数按十进制位折叠，每步都用快速幂并取模。把这个特例继续拆开看：每一位或每个掩码位都要对应题目里的一个布尔事实。位运算只是压缩表示；如果两个状态在位置相同但掩码不同，未来能力不同，就不能合并。

\textbf{状态复用。} 把集合状态压成掩码；包含、加入、删除、枚举子集都变成位操作，但每一位的题目含义不能丢。本题落点是：规律是把大指数按十进制位折叠，每步都用快速幂并取模。手算时只改被新输入影响的那一小部分，而不是回头重算完整候选。

\textbf{指针和字段。} 状态字段要能说清：掩码含义、当前位、目标位集合、状态转移和答案读取；负数或高位参与时要说明整数宽度。手算时按这个协议跟踪：拿 a=2、b=[1, 0] 手算，目标是 \ensuremath{2^{10}} mod 1337，不能把很长指数转成普通整数。扫描十进制指数时，新答案等于旧答案的十次方乘 a 的当前位次方。然后按协议复查：手算时把整数写成二进制，并在每一位旁边标注含义。状态转移只允许改变题目动作允许改变的那些位。

\textbf{代码落点。} 移位表达式加括号；使用无符号或足够宽类型；枚举子集时确认空集和全集是否包含。关键代码行必须能逐句翻译成“旧状态怎样变成新状态”，否则就只是模板。

\textbf{正确性。} 证明从逐位独立或子集归纳开始：被压缩到位里的信息足以决定未来转移。

\textbf{边界实验。} 先用指数为空、a 大于模数、指数含零、递归深度、模乘溢出做反例检查。实验时覆盖指数为空、a 大于模数、指数含零、递归深度、模乘溢出，把数字写成二进制逐位检查。负数、最高位、空掩码和全集掩码都要单独测。

\textbf{复杂度变化。} 暴力用集合或数组保存状态可能需要更多空间和比较成本；位运算通常把每步状态更新压成常数或按位数计算，掩码 DP 仍按状态数增长。

\textbf{迁移追问。} 如果题目改成“若模数变化或要求精确值，取模快速幂的状态语义要重写”，先判断原状态是否仍然覆盖新答案集合，再决定是微调边界、改 value 语义，还是换算法。

\subsection{LeetCode 509 斐波那契数}

\textbf{题目说明。} LeetCode 509 斐波那契数 的题面入口要先还原成具体任务：递归定义会让 fib(n-1) 和 fib(n-2) 下面的状态反复计算。

\textbf{样例入口。} 拿 n=5 手算，fib(5) 会递归调用 fib(4) 和 fib(3)，而 fib(4) 内部又会调用 fib(3)，同一个状态被重复计算。

\textbf{暴力基线。} 慢版本先写递归搜索树：节点表示处理位置、剩余资源和当前选择。只要不同路径反复到达同一个节点，就说明需要把这个节点命名成 DP 状态。放到本题就是：先按“递归定义会让 fib(n-1) 和 fib(n-2) 下面的状态反复计算”完整试慢版本；样例里已经能看到“规律是从 0 和 1 开始向上滚动保存前两项，每个 fib(i) 只算一次”，所以优化前必须先能说清慢版本枚举了哪些候选。

\textbf{暴力过程。} 拿 n=5 手算，fib(5) 会递归调用 fib(4) 和 fib(3)，而 fib(4) 内部又会调用 fib(3)，同一个状态被重复计算。

\textbf{重复工作。} 题面模型：递归定义会让 fib(n-1) 和 fib(n-2) 下面的状态反复计算。暴力版的慢点要落到本题：慢版本会在“递归定义会让 fib(n-1) 和 fib(n-2) 下面的状态反复计算”的选择树里反复到达相同参数。样例规律是“规律是从 0 和 1 开始向上滚动保存前两项，每个 fib(i) 只算一次”，说明这个参数组合可以命名成状态，算过之后后面直接复用。后面的优化只消掉这类重复工作，不能改变答案集合。

\textbf{优化条件。} 本题的关键观察是：自底向上只保存前两项，每次生成当前项后滚动更新。这句话必须能翻译成可维护状态；如果题目条件改变后观察不再成立，就不能继续套原来的写法。

\textbf{相邻状态。} 规律是从 0 和 1 开始向上滚动保存前两项，每个 fib(i) 只算一次。把这个特例继续拆开看：先问暴力搜索里哪些不同路径到达同一个后续问题，再给这个后续问题命名为状态。状态必须包含未来决策需要的全部信息，转移只从更小且已经算清的状态来。

\textbf{状态复用。} 把重复递归节点命名为状态；遍历顺序只是在保证依赖状态已经算完。本题落点是：规律是从 0 和 1 开始向上滚动保存前两项，每个 fib(i) 只算一次。手算时只改被新输入影响的那一小部分，而不是回头重算完整候选。

\textbf{指针和字段。} 状态字段要能说清：状态下标、状态值语义、初始化、转移来源、遍历顺序；空间压缩时还要指出读到的是旧值还是新值。手算时按这个协议跟踪：拿 n=5 手算，fib(5) 会递归调用 fib(4) 和 fib(3)，而 fib(4) 内部又会调用 fib(3)，同一个状态被重复计算。然后按协议复查：手算时先写一个很小的 DP 表或记忆化搜索记录。每个格子旁边写它来自哪些更小状态。

\textbf{代码落点。} 先写未压缩版本校验转移；压缩前后用小表对拍；不可达哨兵参与加法前要检查溢出。关键代码行必须能逐句翻译成“旧状态怎样变成新状态”，否则就只是模板。

\textbf{正确性。} 证明从最后一步开始：任意最优解的最后一步都在转移枚举中，转移来源已经由更小状态正确计算。

\textbf{边界实验。} 先用n 为零、n 为一、结果范围、递归爆栈做反例检查。实验时覆盖n 为零、n 为一、结果范围、递归爆栈，先打印完整 DP 表，再比较空间压缩版本。若压缩版不同，优先检查遍历顺序和旧值保存。

\textbf{复杂度变化。} 暴力递归会重复计算相同子问题，常见指数级；DP 把每个状态算一次，时间是状态数乘单个转移成本，空间是状态表大小或压缩后的大小。

\textbf{迁移追问。} 如果题目改成“若 n 极大，可用矩阵快速幂或快速倍增”，先判断原状态是否仍然覆盖新答案集合，再决定是微调边界、改 value 语义，还是换算法。

\subsection{LeetCode 790 多米诺和托米诺平铺}

\textbf{题目说明。} LeetCode 790 多米诺和托米诺平铺 的题面入口要先还原成具体任务：铺满前 i 列时，边界可能完整，也可能上缺或下缺一个格子。

\textbf{样例入口。} 拿 n=3 手算，只有多米诺可以铺出若干完整状态，加入 L 形托米诺后会产生上缺一格或下缺一格的半填状态。

\textbf{暴力基线。} 慢版本先写递归搜索树：节点表示处理位置、剩余资源和当前选择。只要不同路径反复到达同一个节点，就说明需要把这个节点命名成 DP 状态。放到本题就是：先按“铺满前 i 列时，边界可能完整，也可能上缺或下缺一个格子”完整试慢版本；样例里已经能看到“规律是 DP 不能只保存完整列，还要保存缺口状态；完整状态和两个缺口状态互相转移”，所以优化前必须先能说清慢版本枚举了哪些候选。

\textbf{暴力过程。} 拿 n=3 手算，只有多米诺可以铺出若干完整状态，加入 L 形托米诺后会产生上缺一格或下缺一格的半填状态。

\textbf{重复工作。} 题面模型：铺满前 i 列时，边界可能完整，也可能上缺或下缺一个格子。暴力版的慢点要落到本题：慢版本会在“铺满前 i 列时，边界可能完整，也可能上缺或下缺一个格子”的选择树里反复到达相同参数。样例规律是“规律是 DP 不能只保存完整列，还要保存缺口状态；完整状态和两个缺口状态互相转移”，说明这个参数组合可以命名成状态，算过之后后面直接复用。后面的优化只消掉这类重复工作，不能改变答案集合。

\textbf{优化条件。} 本题的关键观察是：维护完整状态和缺口状态，按最后一块多米诺或托米诺转移。这句话必须能翻译成可维护状态；如果题目条件改变后观察不再成立，就不能继续套原来的写法。

\textbf{相邻状态。} 规律是 DP 不能只保存完整列，还要保存缺口状态；完整状态和两个缺口状态互相转移。把这个特例继续拆开看：先问暴力搜索里哪些不同路径到达同一个后续问题，再给这个后续问题命名为状态。状态必须包含未来决策需要的全部信息，转移只从更小且已经算清的状态来。

\textbf{状态复用。} 把重复递归节点命名为状态；遍历顺序只是在保证依赖状态已经算完。本题落点是：规律是 DP 不能只保存完整列，还要保存缺口状态；完整状态和两个缺口状态互相转移。手算时只改被新输入影响的那一小部分，而不是回头重算完整候选。

\textbf{指针和字段。} 状态字段要能说清：状态下标、状态值语义、初始化、转移来源、遍历顺序；空间压缩时还要指出读到的是旧值还是新值。手算时按这个协议跟踪：拿 n=3 手算，只有多米诺可以铺出若干完整状态，加入 L 形托米诺后会产生上缺一格或下缺一格的半填状态。然后按协议复查：手算时先写一个很小的 DP 表或记忆化搜索记录。每个格子旁边写它来自哪些更小状态。

\textbf{代码落点。} 先写未压缩版本校验转移；压缩前后用小表对拍；不可达哨兵参与加法前要检查溢出。关键代码行必须能逐句翻译成“旧状态怎样变成新状态”，否则就只是模板。

\textbf{正确性。} 证明从最后一步开始：任意最优解的最后一步都在转移枚举中，转移来源已经由更小状态正确计算。

\textbf{边界实验。} 先用n 为一、n 为二、取模、缺口上下对称、初始化做反例检查。实验时覆盖n 为一、n 为二、取模、缺口上下对称、初始化，先打印完整 DP 表，再比较空间压缩版本。若压缩版不同，优先检查遍历顺序和旧值保存。

\textbf{复杂度变化。} 暴力递归会重复计算相同子问题，常见指数级；DP 把每个状态算一次，时间是状态数乘单个转移成本，空间是状态表大小或压缩后的大小。

\textbf{迁移追问。} 如果题目改成“更多形状的铺砖题通常要改成轮廓状态压缩 DP”，先判断原状态是否仍然覆盖新答案集合，再决定是微调边界、改 value 语义，还是换算法。

\subsection{LeetCode 28 找出字符串中第一个匹配项的下标}

\textbf{题目说明。} LeetCode 28 找出字符串中第一个匹配项的下标 的题面入口要先还原成具体任务：字符串匹配的暴力回退会反复比较已经匹配过的前缀。

\textbf{样例入口。} 拿 haystack=sadbutsad、needle=sad 手算，暴力从每个起点尝试匹配，起点 0 成功；若文本和模式有大量重复前缀，失配后从下一个起点重比会浪费。

\textbf{暴力基线。} 暴力从 haystack 的每个起点尝试匹配 needle，遇到失配就把起点加一，needle 从头再比。

\textbf{暴力过程。} 以 haystack=sadbutsad、needle=sad 为例，起点 0 成功。若换成很多重复前缀的文本，如 aaaaaab 匹配 aaaab，暴力会反复比较相同的 a。

\textbf{重复工作。} 题面模型：字符串匹配的暴力回退会反复比较已经匹配过的前缀。暴力版的慢点要落到本题：浪费在失配后完全丢掉已经匹配的前缀信息。模式串内部若有相同前后缀，下一次匹配可以复用一部分。后面的优化只消掉这类重复工作，不能改变答案集合。

\textbf{优化条件。} 本题的关键观察是：KMP 用前缀函数记录模式串最长相等前后缀，失配时把模式串跳到可复用前缀位置。这句话必须能翻译成可维护状态；如果题目条件改变后观察不再成立，就不能继续套原来的写法。

\textbf{相邻状态。} 规律是 KMP 用模式串前缀函数记录可复用的最长相等前后缀，失配时移动模式而不是回退文本。把这个特例继续拆开看：先标出已经确认的前缀或后缀，再标出未知区；能被丢掉的候选，必须是以后再也不会改变已确认区域语义的候选。这样才算从例子推出数组不变量，而不是只记一次操作。

\textbf{状态复用。} KMP 预处理 lps[i]，表示 needle[0..i] 的最长相等真前后缀长度。匹配时文本指针不回退，模式指针按 lps 跳转。手算时只改被新输入影响的那一小部分，而不是回头重算完整候选。

\textbf{指针和字段。} 状态字段要能说清：i 是文本位置，j 是模式位置，lps[j-1] 是失配后可复用的模式前缀长度。手算时按这个协议跟踪：模式 aaaab 在第 5 个字符失配时，前面 aaa 已经可作为下一轮前缀，不必让文本回到旧起点重新比较。

\textbf{代码落点。} 关键代码是失配时 if (j>0) j=lps[j-1]; else ++i。匹配到 j==m 时返回 i-m。关键代码行必须能逐句翻译成“旧状态怎样变成新状态”，否则就只是模板。

\textbf{正确性。} lps 保存的是已经匹配部分中最长可复用前缀，跳转后所有被跳过的起点都无法比这个前缀匹配更多字符。

\textbf{边界实验。} 先用空模式串、模式比文本长、重复字符、完全不匹配、重叠匹配做反例检查。实验时把空模式串、模式比文本长、重复字符、完全不匹配、重叠匹配列成表，再打印关键下标和有效区间。若某个元素被写入后又被未来步骤破坏，说明区间不变量没有成立。

\textbf{复杂度变化。} 暴力最坏 O(nm)，KMP 预处理 O(m)、匹配 O(n)，空间 O(m)。

\textbf{迁移追问。} 如果题目改成“若只做少量短串匹配，暴力可接受；大量匹配再考虑 KMP、哈希或自动机”，先判断原状态是否仍然覆盖新答案集合，再决定是微调边界、改 value 语义，还是换算法。

\subsection{LeetCode 459 重复的子字符串}

\textbf{题目说明。} LeetCode 459 重复的子字符串 的题面入口要先还原成具体任务：判断字符串能否由某个较短周期重复组成，本质是找合法周期长度。

\textbf{样例入口。} 拿 s=abab 手算，最短周期 ab 重复两次；拿 aba 手算，虽然有前后缀 a，但剩余周期长度 2 不能整除 3。

\textbf{暴力基线。} 暴力枚举可能周期 len，若 n 能被 len 整除，就逐段检查 s 是否由 s[0..len-1] 重复组成。

\textbf{暴力过程。} 以 s=abab 为例，len=1 失败，len=2 的周期 ab 重复两次成功；s=aba 虽有前后缀 a，但周期长度 2 不能整除 3。

\textbf{重复工作。} 题面模型：判断字符串能否由某个较短周期重复组成，本质是找合法周期长度。暴力版的慢点要落到本题：浪费在每个周期都重新比较。KMP 前缀函数已经记录了最长相等前后缀，从而推出最小可能周期。后面的优化只消掉这类重复工作，不能改变答案集合。

\textbf{优化条件。} 本题的关键观察是：可用 KMP 前缀函数得到最长相等前后缀，若剩余周期长度能整除总长则成立。这句话必须能翻译成可维护状态；如果题目条件改变后观察不再成立，就不能继续套原来的写法。

\textbf{相邻状态。} 规律是 KMP 前缀函数给出最长相等前后缀，周期长度为 n-lps，能整除才说明由子串重复组成。把这个特例继续拆开看：先标出已经确认的前缀或后缀，再标出未知区；能被丢掉的候选，必须是以后再也不会改变已确认区域语义的候选。这样才算从例子推出数组不变量，而不是只记一次操作。

\textbf{状态复用。} 计算 lps[n-1]，令 period=n-lps[n-1]。若 lps>0 且 n\%period==0，则可以由长度 period 的子串重复。手算时只改被新输入影响的那一小部分，而不是回头重算完整候选。

\textbf{指针和字段。} 状态字段要能说清：lps[i] 是 s[0..i] 的最长相等真前后缀长度，period 是候选重复单元长度。手算时按这个协议跟踪：abab 的 lps 最后为 2，period=2，4\%2==0，所以成立；aba 的 lps 最后为 1，period=2，3\%2!=0，所以不成立。

\textbf{代码落点。} 关键是构建前缀函数时失配用 j=lps[j-1] 回退；最终还要检查整除条件，不能只看 lps>0。关键代码行必须能逐句翻译成“旧状态怎样变成新状态”，否则就只是模板。

\textbf{正确性。} 若字符串由周期重复组成，则存在长度 n-period 的相等前后缀；反过来，最长前后缀推出的 period 能整除 n 时，整串可按该周期铺满。

\textbf{边界实验。} 先用长度一、全相同字符、无重复、部分前后缀相同但不能整除做反例检查。实验时把长度一、全相同字符、无重复、部分前后缀相同但不能整除列成表，再打印关键下标和有效区间。若某个元素被写入后又被未来步骤破坏，说明区间不变量没有成立。

\textbf{复杂度变化。} 枚举周期最坏 O(\ensuremath{n^{2}})，KMP 前缀函数 O(n) 时间、O(n) 空间。

\textbf{迁移追问。} 如果题目改成“也可用 s+s 去掉首尾后查找原串，但要能解释为什么排除了自身匹配”，先判断原状态是否仍然覆盖新答案集合，再决定是微调边界、改 value 语义，还是换算法。

\subsection{LeetCode 208 实现 Trie}

\textbf{题目说明。} LeetCode 208 实现 Trie 的题面入口要先还原成具体任务：前缀树把字符串公共前缀共享成路径。

\textbf{样例入口。} 拿插入 apple 再查 app 手算，app 的路径存在，但如果 p 节点没有终止标记，search(app) 应为假，startsWith(app) 才为真。

\textbf{暴力基线。} 慢版本在每个节点重新扫描子树或路径，先求出局部答案。重复扫描暴露了真正状态：每个节点只需要从孩子或父亲那里拿到少量信息。放到本题就是：先按“前缀树把字符串公共前缀共享成路径”完整试慢版本；样例里已经能看到“规律是 Trie 节点既要有孩子边，也要有单词终止标记”，所以优化前必须先能说清慢版本枚举了哪些候选。

\textbf{暴力过程。} 拿插入 apple 再查 app 手算，app 的路径存在，但如果 p 节点没有终止标记，search(app) 应为假，startsWith(app) 才为真。

\textbf{重复工作。} 题面模型：前缀树把字符串公共前缀共享成路径。暴力版的慢点要落到本题：慢版本会在树上重复确认“前缀树把字符串公共前缀共享成路径”。样例规律是“规律是 Trie 节点既要有孩子边，也要有单词终止标记”，说明父子之间真正需要传递的是高度、上下界、命中状态、层号或两类选择值，而不是反复扫描整棵子树。后面的优化只消掉这类重复工作，不能改变答案集合。

\textbf{优化条件。} 本题的关键观察是：插入和查询都逐字符沿边移动，不存在的边按需要创建。这句话必须能翻译成可维护状态；如果题目条件改变后观察不再成立，就不能继续套原来的写法。

\textbf{相邻状态。} 规律是 Trie 节点既要有孩子边，也要有单词终止标记。把这个特例继续拆开看：节点向父亲返回的量和全局答案通常不是同一个量。先写叶子或空节点的返回值，再写父节点如何合并左右孩子，递归式才有边界基础。

\textbf{状态复用。} 把对子树的重复扫描压成返回值或队列状态；每个节点只合并孩子给出的稳定信息。本题落点是：规律是 Trie 节点既要有孩子边，也要有单词终止标记。手算时只改被新输入影响的那一小部分，而不是回头重算完整候选。

\textbf{指针和字段。} 状态字段要能说清：节点指针、传入约束或返回值、当前答案、层号或路径状态；空节点的返回值必须和状态定义匹配。手算时按这个协议跟踪：拿插入 apple 再查 app 手算，app 的路径存在，但如果 p 节点没有终止标记，search(app) 应为假，startsWith(app) 才为真。然后按协议复查：手算时从叶子开始写返回值，或者从根开始写传入约束。不要只写访问顺序，要写每条边上传递的信息。

\textbf{代码落点。} 递归版本先处理空节点；迭代版本的栈元素要带完整状态；全负值、重复值和极端深树要单独测。关键代码行必须能逐句翻译成“旧状态怎样变成新状态”，否则就只是模板。

\textbf{正确性。} 证明从子树归纳或层序不变量开始：孩子状态正确时父节点合并正确；若用队列，当前轮队列必须正好保存当前层节点。

\textbf{边界实验。} 先用空字符串约定、单词结束标记、前缀存在不等于单词存在做反例检查。实验时覆盖空字符串约定、单词结束标记、前缀存在不等于单词存在，尤其关注空节点、单节点、链式树和全负值。递归版本和显式栈版本可以互相对拍。

\textbf{复杂度变化。} 暴力在每个节点重复扫描子树或反复寻找层边界可能退化到平方级；后序汇总、显式栈或层序遍历让每个节点处理常数次，通常是线性时间。

\textbf{迁移追问。} 如果题目改成“若字符集很大，用哈希子节点；小写字母用数组更快”，先判断原状态是否仍然覆盖新答案集合，再决定是微调边界、改 value 语义，还是换算法。

\subsection{LeetCode 211 添加与搜索单词}

\textbf{题目说明。} LeetCode 211 添加与搜索单词 的题面入口要先还原成具体任务：点号通配符让查询可能分叉，但前缀树仍能剪掉不匹配前缀。

\textbf{样例入口。} 拿 add bad、dad、mad 后查 .ad 手算，点号可以走 b、d、m 三条边；查 b.. 时后两个点继续在子树里展开。

\textbf{暴力基线。} 慢版本在每个节点重新扫描子树或路径，先求出局部答案。重复扫描暴露了真正状态：每个节点只需要从孩子或父亲那里拿到少量信息。放到本题就是：先按“点号通配符让查询可能分叉，但前缀树仍能剪掉不匹配前缀”完整试慢版本；样例里已经能看到“规律是普通字符沿唯一边走，点号触发 DFS 枚举孩子，终点仍要检查单词结束标记”，所以优化前必须先能说清慢版本枚举了哪些候选。

\textbf{暴力过程。} 拿 add bad、dad、mad 后查 .ad 手算，点号可以走 b、d、m 三条边；查 b.. 时后两个点继续在子树里展开。

\textbf{重复工作。} 题面模型：点号通配符让查询可能分叉，但前缀树仍能剪掉不匹配前缀。暴力版的慢点要落到本题：慢版本会在树上重复确认“点号通配符让查询可能分叉，但前缀树仍能剪掉不匹配前缀”。样例规律是“规律是普通字符沿唯一边走，点号触发 DFS 枚举孩子，终点仍要检查单词结束标记”，说明父子之间真正需要传递的是高度、上下界、命中状态、层号或两类选择值，而不是反复扫描整棵子树。后面的优化只消掉这类重复工作，不能改变答案集合。

\textbf{优化条件。} 本题的关键观察是：普通字符沿唯一边走，点号枚举当前节点所有孩子。这句话必须能翻译成可维护状态；如果题目条件改变后观察不再成立，就不能继续套原来的写法。

\textbf{相邻状态。} 规律是普通字符沿唯一边走，点号触发 DFS 枚举孩子，终点仍要检查单词结束标记。把这个特例继续拆开看：节点向父亲返回的量和全局答案通常不是同一个量。先写叶子或空节点的返回值，再写父节点如何合并左右孩子，递归式才有边界基础。

\textbf{状态复用。} 把对子树的重复扫描压成返回值或队列状态；每个节点只合并孩子给出的稳定信息。本题落点是：规律是普通字符沿唯一边走，点号触发 DFS 枚举孩子，终点仍要检查单词结束标记。手算时只改被新输入影响的那一小部分，而不是回头重算完整候选。

\textbf{指针和字段。} 状态字段要能说清：节点指针、传入约束或返回值、当前答案、层号或路径状态；空节点的返回值必须和状态定义匹配。手算时按这个协议跟踪：拿 add bad、dad、mad 后查 .ad 手算，点号可以走 b、d、m 三条边；查 b.. 时后两个点继续在子树里展开。然后按协议复查：手算时从叶子开始写返回值，或者从根开始写传入约束。不要只写访问顺序，要写每条边上传递的信息。

\textbf{代码落点。} 递归版本先处理空节点；迭代版本的栈元素要带完整状态；全负值、重复值和极端深树要单独测。关键代码行必须能逐句翻译成“旧状态怎样变成新状态”，否则就只是模板。

\textbf{正确性。} 证明从子树归纳或层序不变量开始：孩子状态正确时父节点合并正确；若用队列，当前轮队列必须正好保存当前层节点。

\textbf{边界实验。} 先用连续通配符、空结果、单词结束标记、字符集大小做反例检查。实验时覆盖连续通配符、空结果、单词结束标记、字符集大小，尤其关注空节点、单节点、链式树和全负值。递归版本和显式栈版本可以互相对拍。

\textbf{复杂度变化。} 暴力在每个节点重复扫描子树或反复寻找层边界可能退化到平方级；后序汇总、显式栈或层序遍历让每个节点处理常数次，通常是线性时间。

\textbf{迁移追问。} 如果题目改成“大量通配符会接近遍历整棵 Trie，需要规模约束”，先判断原状态是否仍然覆盖新答案集合，再决定是微调边界、改 value 语义，还是换算法。

\subsection{LeetCode 785 判断二分图}

\textbf{题目说明。} LeetCode 785 判断二分图 的题面入口要先还原成具体任务：二分图要求每条边两端颜色相反。

\textbf{样例入口。} 拿 graph=[[1, 2, 3], [0, 2], [0, 1, 3], [0, 2]] 手算，0 染红后 1、2、3 都要染蓝，但 1 和 2 之间还有边，蓝蓝相邻冲突，所以不是二分图。

\textbf{暴力基线。} 慢版本先按题意画出完整状态图，用普通搜索跑通可达性。这个过程用于确认二分图要求每条边两端颜色相反的节点和边，之后再讨论访问标记、层次或代价顺序。放到本题就是：先按“二分图要求每条边两端颜色相反”完整试慢版本；样例里已经能看到“规律是 BFS 染色时每条边两端必须异色，发现同色边立即失败”，所以优化前必须先能说清慢版本枚举了哪些候选。

\textbf{暴力过程。} 拿 graph=[[1, 2, 3], [0, 2], [0, 1, 3], [0, 2]] 手算，0 染红后 1、2、3 都要染蓝，但 1 和 2 之间还有边，蓝蓝相邻冲突，所以不是二分图。

\textbf{重复工作。} 题面模型：二分图要求每条边两端颜色相反。暴力版的慢点要落到本题：慢版本会围绕“二分图要求每条边两端颜色相反”反复展开同一批状态。样例规律是“规律是 BFS 染色时每条边两端必须异色，发现同色边立即失败”，说明必须把节点、边、访问标记和资源维度一次建清；同一状态不应重复入队，不同资源状态也不能误合并。后面的优化只消掉这类重复工作，不能改变答案集合。

\textbf{优化条件。} 本题的关键观察是：对每个未染色连通块 BFS 或 DFS 染色，遇到同色边立即失败。这句话必须能翻译成可维护状态；如果题目条件改变后观察不再成立，就不能继续套原来的写法。

\textbf{相邻状态。} 规律是 BFS 染色时每条边两端必须异色，发现同色边立即失败。把这个特例继续拆开看：状态节点不一定等于原始位置，还可能包含钥匙、剩余资源、时间或访问集合。visited 只能合并未来能力完全相同的状态，否则会把合法路径剪掉。

\textbf{状态复用。} 把重复搜索压成访问标记、距离数组或拓扑状态；邻居生成也要从全量扫描变成结构化枚举。本题落点是：规律是 BFS 染色时每条边两端必须异色，发现同色边立即失败。手算时只改被新输入影响的那一小部分，而不是回头重算完整候选。

\textbf{指针和字段。} 状态字段要能说清：状态编号、邻接生成方式、访问标记、距离或层数、队列内容；若状态含资源，visited 不能只按位置记录。手算时按这个协议跟踪：拿 graph=[[1, 2, 3], [0, 2], [0, 1, 3], [0, 2]] 手算，0 染红后 1、2、3 都要染蓝，但 1 和 2 之间还有边，蓝蓝相邻冲突，所以不是二分图。然后按协议复查：手算时画队列或栈，状态必须包含题目需要的所有资源。入队时标记还是出队时标记，要和去重语义一致。

\textbf{代码落点。} 入队时标记还是出队时标记必须一致；方向数组集中定义；多源 BFS 要一次性把所有源点放入初始队列。关键代码行必须能逐句翻译成“旧状态怎样变成新状态”，否则就只是模板。

\textbf{正确性。} 证明从可达性开始：搜索覆盖所有合法边能到达的状态，访问标记只删除重复状态而不删除不同未来能力。

\textbf{边界实验。} 先用非连通图、自环、奇环、孤立点、重复边做反例检查。实验时覆盖非连通图、自环、奇环、孤立点、重复边，并打印入队时的完整状态。若同一位置但资源不同的状态被合并，很多图题会出现隐蔽错误。

\textbf{复杂度变化。} 暴力重复从多个起点搜索会反复遍历图；正确建模和访问标记后，BFS 或 DFS 通常按节点数加边数计算，状态图题则按完整状态数计算。

\textbf{迁移追问。} 如果题目改成“可能的二分法是同一染色模型，只是图由 dislike 关系构造”，先判断原状态是否仍然覆盖新答案集合，再决定是微调边界、改 value 语义，还是换算法。

\subsection{LeetCode 886 可能的二分法}

\textbf{题目说明。} LeetCode 886 可能的二分法 的题面入口要先还原成具体任务：互相讨厌的人必须分到不同组，本质是图二染色。

\textbf{样例入口。} 拿 dislikes=[[1, 2], [1, 3], [2, 4]] 手算，把 1 放 A，则 2 和 3 必须放 B，2 讨厌 4，所以 4 放 A，可行。若出现奇环，某条边两端会被迫同色。

\textbf{暴力基线。} 慢版本先按题意画出完整状态图，用普通搜索跑通可达性。这个过程用于确认互相讨厌的人必须分到不同组，本质是图二染色的节点和边，之后再讨论访问标记、层次或代价顺序。放到本题就是：先按“互相讨厌的人必须分到不同组，本质是图二染色”完整试慢版本；样例里已经能看到“规律是把 dislike 当无向边做二染色，发现同色冲突就不能二分”，所以优化前必须先能说清慢版本枚举了哪些候选。

\textbf{暴力过程。} 拿 dislikes=[[1, 2], [1, 3], [2, 4]] 手算，把 1 放 A，则 2 和 3 必须放 B，2 讨厌 4，所以 4 放 A，可行。若出现奇环，某条边两端会被迫同色。

\textbf{重复工作。} 题面模型：互相讨厌的人必须分到不同组，本质是图二染色。暴力版的慢点要落到本题：慢版本会围绕“互相讨厌的人必须分到不同组，本质是图二染色”反复展开同一批状态。样例规律是“规律是把 dislike 当无向边做二染色，发现同色冲突就不能二分”，说明必须把节点、边、访问标记和资源维度一次建清；同一状态不应重复入队，不同资源状态也不能误合并。后面的优化只消掉这类重复工作，不能改变答案集合。

\textbf{优化条件。} 本题的关键观察是：按 dislike 建无向图，对每个连通块 BFS 染色，发现同色冲突则不可行。这句话必须能翻译成可维护状态；如果题目条件改变后观察不再成立，就不能继续套原来的写法。

\textbf{相邻状态。} 规律是把 dislike 当无向边做二染色，发现同色冲突就不能二分。把这个特例继续拆开看：状态节点不一定等于原始位置，还可能包含钥匙、剩余资源、时间或访问集合。visited 只能合并未来能力完全相同的状态，否则会把合法路径剪掉。

\textbf{状态复用。} 把重复搜索压成访问标记、距离数组或拓扑状态；邻居生成也要从全量扫描变成结构化枚举。本题落点是：规律是把 dislike 当无向边做二染色，发现同色冲突就不能二分。手算时只改被新输入影响的那一小部分，而不是回头重算完整候选。

\textbf{指针和字段。} 状态字段要能说清：状态编号、邻接生成方式、访问标记、距离或层数、队列内容；若状态含资源，visited 不能只按位置记录。手算时按这个协议跟踪：拿 dislikes=[[1, 2], [1, 3], [2, 4]] 手算，把 1 放 A，则 2 和 3 必须放 B，2 讨厌 4，所以 4 放 A，可行。若出现奇环，某条边两端会被迫同色。然后按协议复查：手算时画队列或栈，状态必须包含题目需要的所有资源。入队时标记还是出队时标记，要和去重语义一致。

\textbf{代码落点。} 入队时标记还是出队时标记必须一致；方向数组集中定义；多源 BFS 要一次性把所有源点放入初始队列。关键代码行必须能逐句翻译成“旧状态怎样变成新状态”，否则就只是模板。

\textbf{正确性。} 证明从可达性开始：搜索覆盖所有合法边能到达的状态，访问标记只删除重复状态而不删除不同未来能力。

\textbf{边界实验。} 先用孤立点、非连通图、奇环、重复 dislike、自身 dislike做反例检查。实验时覆盖孤立点、非连通图、奇环、重复 dislike、自身 dislike，并打印入队时的完整状态。若同一位置但资源不同的状态被合并，很多图题会出现隐蔽错误。

\textbf{复杂度变化。} 暴力重复从多个起点搜索会反复遍历图；正确建模和访问标记后，BFS 或 DFS 通常按节点数加边数计算，状态图题则按完整状态数计算。

\textbf{迁移追问。} 如果题目改成“它和判断二分图同型，只是输入不是邻接表”，先判断原状态是否仍然覆盖新答案集合，再决定是微调边界、改 value 语义，还是换算法。

\subsection{LeetCode 600 不含连续 1 的非负整数}

\textbf{题目说明。} LeetCode 600 不含连续 1 的非负整数 的题面入口要先还原成具体任务：统计不超过 n 的二进制数，需要按位构造并记住上一位是否为一。

\textbf{样例入口。} 拿 n=5 的二进制 101 手算，不含连续 1 的数是 0、1、2、4、5。按高位扫描时，遇到某位为 1，可以先把这一位改成 0，后面任意填不含连续 1 的低位；若扫描中出现连续 1，就不能继续把 n 本身计入。

\textbf{暴力基线。} 慢版本先写递归搜索树：节点表示处理位置、剩余资源和当前选择。只要不同路径反复到达同一个节点，就说明需要把这个节点命名成 DP 状态。放到本题就是：先按“统计不超过 n 的二进制数，需要按位构造并记住上一位是否为一”完整试慢版本；样例里已经能看到“规律是数位 DP 需要保存上一位是否为 1 和是否贴着上界”，所以优化前必须先能说清慢版本枚举了哪些候选。

\textbf{暴力过程。} 拿 n=5 的二进制 101 手算，不含连续 1 的数是 0、1、2、4、5。按高位扫描时，遇到某位为 1，可以先把这一位改成 0，后面任意填不含连续 1 的低位；若扫描中出现连续 1，就不能继续把 n 本身计入。

\textbf{重复工作。} 题面模型：统计不超过 n 的二进制数，需要按位构造并记住上一位是否为一。暴力版的慢点要落到本题：慢版本会在“统计不超过 n 的二进制数，需要按位构造并记住上一位是否为一”的选择树里反复到达相同参数。样例规律是“规律是数位 DP 需要保存上一位是否为 1 和是否贴着上界”，说明这个参数组合可以命名成状态，算过之后后面直接复用。后面的优化只消掉这类重复工作，不能改变答案集合。

\textbf{优化条件。} 本题的关键观察是：数位 DP 或 Fibonacci 计数从高位扫描，遇到 1 时累加较小前缀数量，连续 1 时停止。这句话必须能翻译成可维护状态；如果题目条件改变后观察不再成立，就不能继续套原来的写法。

\textbf{相邻状态。} 规律是数位 DP 需要保存上一位是否为 1 和是否贴着上界。把这个特例继续拆开看：先问暴力搜索里哪些不同路径到达同一个后续问题，再给这个后续问题命名为状态。状态必须包含未来决策需要的全部信息，转移只从更小且已经算清的状态来。

\textbf{状态复用。} 把重复递归节点命名为状态；遍历顺序只是在保证依赖状态已经算完。本题落点是：规律是数位 DP 需要保存上一位是否为 1 和是否贴着上界。手算时只改被新输入影响的那一小部分，而不是回头重算完整候选。

\textbf{指针和字段。} 状态字段要能说清：状态下标、状态值语义、初始化、转移来源、遍历顺序；空间压缩时还要指出读到的是旧值还是新值。手算时按这个协议跟踪：拿 n=5 的二进制 101 手算，不含连续 1 的数是 0、1、2、4、5。按高位扫描时，遇到某位为 1，可以先把这一位改成 0，后面任意填不含连续 1 的低位；若扫描中出现连续 1，就不能继续把 n 本身计入。然后按协议复查：手算时先写一个很小的 DP 表或记忆化搜索记录。每个格子旁边写它来自哪些更小状态。

\textbf{代码落点。} 先写未压缩版本校验转移；压缩前后用小表对拍；不可达哨兵参与加法前要检查溢出。关键代码行必须能逐句翻译成“旧状态怎样变成新状态”，否则就只是模板。

\textbf{正确性。} 证明从最后一步开始：任意最优解的最后一步都在转移枚举中，转移来源已经由更小状态正确计算。

\textbf{边界实验。} 先用n 为零、最高位、连续一提前停止、包含 n 本身、整数位宽做反例检查。实验时覆盖n 为零、最高位、连续一提前停止、包含 n 本身、整数位宽，先打印完整 DP 表，再比较空间压缩版本。若压缩版不同，优先检查遍历顺序和旧值保存。

\textbf{复杂度变化。} 暴力递归会重复计算相同子问题，常见指数级；DP 把每个状态算一次，时间是状态数乘单个转移成本，空间是状态表大小或压缩后的大小。

\textbf{迁移追问。} 如果题目改成“若禁止其他数字模式，状态要记录更长后缀或自动机”，先判断原状态是否仍然覆盖新答案集合，再决定是微调边界、改 value 语义，还是换算法。

\subsection{LeetCode 1755 最接近目标值的子序列和}

\textbf{题目说明。} LeetCode 1755 最接近目标值的子序列和 的题面入口要先还原成具体任务：完整子序列选择是 \ensuremath{2^{n}}，拆成左右两半后各自枚举子集和。

\textbf{样例入口。} 拿 nums=[5, -7, 3, 5]、goal=6 手算，完整枚举 16 个子序列可找到 -7+3+5+5=6。折半后左半 [5, -7] 生成 0、5、-7、-2，右半 [3, 5] 生成 0、3、5、8；对每个左半和二分找最接近 goal-left 的右半和。

\textbf{暴力基线。} 慢版本先用集合、数组或布尔表保存状态，逐步执行题目动作。确认每个布尔字段的含义后，才能把它压缩成位，而不是直接写位运算公式。放到本题就是：先按“完整子序列选择是 \ensuremath{2^{n}}，拆成左右两半后各自枚举子集和”完整试慢版本；样例里已经能看到“规律是把 \ensuremath{2^{n}} 拆成两个 \ensuremath{2^{n/2}} 的子集和集合再配对”，所以优化前必须先能说清慢版本枚举了哪些候选。

\textbf{暴力过程。} 拿 nums=[5, -7, 3, 5]、goal=6 手算，完整枚举 16 个子序列可找到 -7+3+5+5=6。折半后左半 [5, -7] 生成 0、5、-7、-2，右半 [3, 5] 生成 0、3、5、8；对每个左半和二分找最接近 goal-left 的右半和。

\textbf{重复工作。} 题面模型：完整子序列选择是 \ensuremath{2^{n}}，拆成左右两半后各自枚举子集和。暴力版的慢点要落到本题：慢版本会围绕“完整子序列选择是 \ensuremath{2^{n}}，拆成左右两半后各自枚举子集和”使用集合、计数或完整状态表。样例规律是“规律是把 \ensuremath{2^{n}} 拆成两个 \ensuremath{2^{n/2}} 的子集和集合再配对”，说明哪些信息可以被逐位抵消、压缩或复用；位运算只是编码，不是证明本身。后面的优化只消掉这类重复工作，不能改变答案集合。

\textbf{优化条件。} 本题的关键观察是：分别生成两半所有子集和，排序一侧，对另一侧每个和用二分找最接近 goal 的补数。这句话必须能翻译成可维护状态；如果题目条件改变后观察不再成立，就不能继续套原来的写法。

\textbf{相邻状态。} 规律是把 \ensuremath{2^{n}} 拆成两个 \ensuremath{2^{n/2}} 的子集和集合再配对。把这个特例继续拆开看：每一位或每个掩码位都要对应题目里的一个布尔事实。位运算只是压缩表示；如果两个状态在位置相同但掩码不同，未来能力不同，就不能合并。

\textbf{状态复用。} 把集合状态压成掩码；包含、加入、删除、枚举子集都变成位操作，但每一位的题目含义不能丢。本题落点是：规律是把 \ensuremath{2^{n}} 拆成两个 \ensuremath{2^{n/2}} 的子集和集合再配对。手算时只改被新输入影响的那一小部分，而不是回头重算完整候选。

\textbf{指针和字段。} 状态字段要能说清：掩码含义、当前位、目标位集合、状态转移和答案读取；负数或高位参与时要说明整数宽度。手算时按这个协议跟踪：拿 nums=[5, -7, 3, 5]、goal=6 手算，完整枚举 16 个子序列可找到 -7+3+5+5=6。折半后左半 [5, -7] 生成 0、5、-7、-2，右半 [3, 5] 生成 0、3、5、8；对每个左半和二分找最接近 goal-left 的右半和。然后按协议复查：手算时把整数写成二进制，并在每一位旁边标注含义。状态转移只允许改变题目动作允许改变的那些位。

\textbf{代码落点。} 移位表达式加括号；使用无符号或足够宽类型；枚举子集时确认空集和全集是否包含。关键代码行必须能逐句翻译成“旧状态怎样变成新状态”，否则就只是模板。

\textbf{正确性。} 证明从逐位独立或子集归纳开始：被压缩到位里的信息足以决定未来转移。

\textbf{边界实验。} 先用负数、空子序列、目标为负、答案为零提前返回、两半大小做反例检查。实验时覆盖负数、空子序列、目标为负、答案为零提前返回、两半大小，把数字写成二进制逐位检查。负数、最高位、空掩码和全集掩码都要单独测。

\textbf{复杂度变化。} 暴力用集合或数组保存状态可能需要更多空间和比较成本；位运算通常把每步状态更新压成常数或按位数计算，掩码 DP 仍按状态数增长。

\textbf{迁移追问。} 如果题目改成“若要求方案数而非最接近差值，要改成计数配对”，先判断原状态是否仍然覆盖新答案集合，再决定是微调边界、改 value 语义，还是换算法。

\subsection{LeetCode 587 安装栅栏}

\textbf{题目说明。} LeetCode 587 安装栅栏 的题面入口要先还原成具体任务：外层点构成凸包，内部点不会影响围栏边界。

\textbf{样例入口。} 拿一组点里所有边界点围成栅栏手算，内部点不会影响围栏；排序后构造下凸壳时，如果新点让边界向右转，就说明中间点在内部，要弹出。

\textbf{暴力基线。} 暴力可以检查每两个点构成的直线，若所有点都在同一侧，则这条线可能是凸包边。

\textbf{暴力过程。} 若有大量点，暴力要枚举点对并检查所有其他点。排序后从左到右构造边界，可以逐步排除内部点。

\textbf{重复工作。} 题面模型：外层点构成凸包，内部点不会影响围栏边界。暴力版的慢点要落到本题：浪费在反复判断同一批点是否在边界内。单调链维护当前凸壳，遇到右转时中间点已经不可能留在外边界。后面的优化只消掉这类重复工作，不能改变答案集合。

\textbf{优化条件。} 本题的关键观察是：排序后用单调链维护上下凸壳，遇到会让边界右转的点就弹出，题目要求保留共线边界点。这句话必须能翻译成可维护状态；如果题目条件改变后观察不再成立，就不能继续套原来的写法。

\textbf{相邻状态。} 规律是单调链用叉积维护凸壳，题目要求保留共线边界点时等号不能随便弹。把这个特例继续拆开看：先标出已经确认的前缀或后缀，再标出未知区；能被丢掉的候选，必须是以后再也不会改变已确认区域语义的候选。这样才算从例子推出数组不变量，而不是只记一次操作。

\textbf{状态复用。} 按坐标排序，分别构造下凸壳和上凸壳。用叉积判断最近两个壳点加新点是否破坏外边界；题目要求保留共线边界点时等号处理要谨慎。手算时只改被新输入影响的那一小部分，而不是回头重算完整候选。

\textbf{指针和字段。} 状态字段要能说清：hull 是当前边界点栈，cross(a, b, c) 表示转向，排序后的点序提供扫描方向。手算时按这个协议跟踪：新点让最后两条边形成右转时，倒数第二个点落到内部，弹出；若三点共线且在边界上，需要保留共线点。

\textbf{代码落点。} 关键代码是 while hull size>=2 \&\& cross(...) < 0 才弹出；若写成 <=0 会把共线边界点删掉。关键代码行必须能逐句翻译成“旧状态怎样变成新状态”，否则就只是模板。

\textbf{正确性。} 单调链每次弹出的点都位于新边形成的内部侧，不可能成为最终外边界；上下壳合并覆盖整个凸包边界。

\textbf{边界实验。} 先用所有点共线、重复点、共线边界是否保留、叉积溢出做反例检查。实验时把所有点共线、重复点、共线边界是否保留、叉积溢出列成表，再打印关键下标和有效区间。若某个元素被写入后又被未来步骤破坏，说明区间不变量没有成立。

\textbf{复杂度变化。} 暴力 O(\ensuremath{n^{3}})，单调链排序 O(n log n)、扫描 O(n)，空间 O(n)。

\textbf{迁移追问。} 如果题目改成“若只求凸包顶点不保留边上点，叉积等号处理不同”，先判断原状态是否仍然覆盖新答案集合，再决定是微调边界、改 value 语义，还是换算法。

\subsection{LeetCode 699 掉落的方块}

\textbf{题目说明。} LeetCode 699 掉落的方块 的题面入口要先还原成具体任务：每个方块会查询一段横坐标当前最高高度，再把该段整体抬高。

\textbf{样例入口。} 拿 positions=[[1, 2], [2, 3], [6, 1]] 手算，第一个方块落在 [1, 3) 高度 2；第二个 [2, 5) 和它重叠，落后高度变 5；第三个不重叠，高度仍取全局最大 5。

\textbf{暴力基线。} 慢版本对新区间和所有旧区间两两比较，手工判断相交、覆盖或冲突。它能验证端点语义，但排序后很多远离当前边界的区间无需再看。放到本题就是：先按“每个方块会查询一段横坐标当前最高高度，再把该段整体抬高”完整试慢版本；样例里已经能看到“规律是每次查询区间当前最大高度，再把该区间赋值为最大高度加方块边长”，所以优化前必须先能说清慢版本枚举了哪些候选。

\textbf{暴力过程。} 拿 positions=[[1, 2], [2, 3], [6, 1]] 手算，第一个方块落在 [1, 3) 高度 2；第二个 [2, 5) 和它重叠，落后高度变 5；第三个不重叠，高度仍取全局最大 5。

\textbf{重复工作。} 题面模型：每个方块会查询一段横坐标当前最高高度，再把该段整体抬高。暴力版的慢点要落到本题：慢版本会围绕“每个方块会查询一段横坐标当前最高高度，再把该段整体抬高”做两两区间比较。样例规律是“规律是每次查询区间当前最大高度，再把该区间赋值为最大高度加方块边长”，说明排序后只有贴近当前端点、结果尾部或资源堆顶的区间还会影响答案。后面的优化只消掉这类重复工作，不能改变答案集合。

\textbf{优化条件。} 本题的关键观察是：坐标压缩后用线段树支持区间最大值查询和区间赋值更新。这句话必须能翻译成可维护状态；如果题目条件改变后观察不再成立，就不能继续套原来的写法。

\textbf{相邻状态。} 规律是每次查询区间当前最大高度，再把该区间赋值为最大高度加方块边长。把这个特例继续拆开看：排序以后，真正还能影响当前区间的候选必须贴近当前端点、结果尾部或资源堆顶。某个区间一旦被覆盖、合并或弹出，就要说明它为什么不会和后续区间重新产生更优关系。

\textbf{状态复用。} 把两两比较压成排序后的相邻关系、当前边界或动态有序结构；远离当前边界的区间要被安全归档。本题落点是：规律是每次查询区间当前最大高度，再把该区间赋值为最大高度加方块边长。手算时只改被新输入影响的那一小部分，而不是回头重算完整候选。

\textbf{指针和字段。} 状态字段要能说清：排序后的区间、当前结果尾、扫描右端或资源堆、端点比较规则；动态版本还要保存前驱后继。手算时按这个协议跟踪：拿 positions=[[1, 2], [2, 3], [6, 1]] 手算，第一个方块落在 [1, 3) 高度 2；第二个 [2, 5) 和它重叠，落后高度变 5；第三个不重叠，高度仍取全局最大 5。然后按协议复查：手算时先排序，再逐个区间写当前结果尾部、当前右端或资源堆。每个区间离开视野时都要说明它不会再影响后续。

\textbf{代码落点。} 排序比较器满足严格弱序；端点相等的处理和题意一致；扫描时只和当前结果尾或当前资源集合交互。关键代码行必须能逐句翻译成“旧状态怎样变成新状态”，否则就只是模板。

\textbf{正确性。} 证明从排序后的邻接关系开始：已经合并、选择或丢弃的区间不会被更后面的区间重新影响。

\textbf{边界实验。} 先用坐标很大、区间端点半开、方块不重叠、完全覆盖、懒标记做反例检查。实验时覆盖坐标很大、区间端点半开、方块不重叠、完全覆盖、懒标记，画出排序后的区间序列和当前结果尾部。动态版本要专门测前驱、后继和端点相接。

\textbf{复杂度变化。} 暴力两两比较区间常见为平方级；排序后扫描通常由排序成本主导，动态有序结构则把每次前驱后继查询控制在对数级。

\textbf{迁移追问。} 如果题目改成“若只离线统计最终高度且无覆盖查询，可以换成扫描线模型”，先判断原状态是否仍然覆盖新答案集合，再决定是微调边界、改 value 语义，还是换算法。

\subsection{LeetCode 778 水位上升的泳池中游泳}

\textbf{题目说明。} LeetCode 778 水位上升的泳池中游泳 的题面入口要先还原成具体任务：到达某格子的代价是路径上最大高度，而不是高度和。

\textbf{样例入口。} 拿 [[0, 2], [1, 3]] 手算，时间必须至少到 3 才能进入右下角；路径代价不是边权和，而是沿路最大高度。

\textbf{暴力基线。} 慢版本可以枚举路径，或先用普通队列试跑一个小图。它会暴露步数最少和代价最小是否一致；一旦不一致，就必须围绕代价组织状态。放到本题就是：先按“到达某格子的代价是路径上最大高度，而不是高度和”完整试慢版本；样例里已经能看到“规律是 Dijkstra 的距离定义为到达某格所需的最小最大高度，或二分时间做可达性”，所以优化前必须先能说清慢版本枚举了哪些候选。

\textbf{暴力过程。} 拿 [[0, 2], [1, 3]] 手算，时间必须至少到 3 才能进入右下角；路径代价不是边权和，而是沿路最大高度。

\textbf{重复工作。} 题面模型：到达某格子的代价是路径上最大高度，而不是高度和。暴力版的慢点要落到本题：慢版本会围绕“到达某格子的代价是路径上最大高度，而不是高度和”枚举路径或按错误顺序扩展状态。样例规律是“规律是 Dijkstra 的距离定义为到达某格所需的最小最大高度，或二分时间做可达性”，说明队列或堆的弹出顺序必须和代价定义一致，旧状态为什么能跳过也要讲清。后面的优化只消掉这类重复工作，不能改变答案集合。

\textbf{优化条件。} 本题的关键观察是：用 Dijkstra 按当前路径最大高度扩展，松弛时取 max 当前代价和邻格高度。这句话必须能翻译成可维护状态；如果题目条件改变后观察不再成立，就不能继续套原来的写法。

\textbf{相邻状态。} 规律是 Dijkstra 的距离定义为到达某格所需的最小最大高度，或二分时间做可达性。把这个特例继续拆开看：队列或堆弹出的顺序必须和代价规则一致；旧状态被跳过，是因为已经有更小代价到达同一状态。只要边权或代价合成方式改变，就要重新证明扩展顺序。

\textbf{状态复用。} 把路径枚举压成距离数组和优先队列；每次只扩展当前最有希望的未过期状态。本题落点是：规律是 Dijkstra 的距离定义为到达某格所需的最小最大高度，或二分时间做可达性。手算时只改被新输入影响的那一小部分，而不是回头重算完整候选。

\textbf{指针和字段。} 状态字段要能说清：距离数组、优先队列状态、松弛候选、旧状态判定、不可达哨兵；资源约束题还要把资源维度放进状态。手算时按这个协议跟踪：拿 [[0, 2], [1, 3]] 手算，时间必须至少到 3 才能进入右下角；路径代价不是边权和，而是沿路最大高度。然后按协议复查：手算时记录每次弹出的代价、当前距离数组和松弛结果。旧状态被弹出时为什么跳过，也要写在表里。

\textbf{代码落点。} 松弛时只在候选更优时更新；priority\_queue 弹出旧状态要跳过；距离加法用更宽整数。关键代码行必须能逐句翻译成“旧状态怎样变成新状态”，否则就只是模板。

\textbf{正确性。} 证明从扩展顺序开始：当前弹出的状态在代价规则下已经不会被未来路径改得更优。

\textbf{边界实验。} 先用单格、起点高度很高、需要绕路、多个相同高度做反例检查。实验时覆盖单格、起点高度很高、需要绕路、多个相同高度，打印每次弹出的代价和是否过期。若普通 BFS 与 Dijkstra 结果不同，要能解释边权条件造成的差异。

\textbf{复杂度变化。} 暴力枚举路径可能指数级；Dijkstra 在邻接表和堆实现下按边数、点数和堆操作计算，0-1 BFS 可按节点数加边数计算。

\textbf{迁移追问。} 如果题目改成“也可以按水位二分可达，或按高度排序并查集合并”，先判断原状态是否仍然覆盖新答案集合，再决定是微调边界、改 value 语义，还是换算法。

\subsection{LeetCode 236 二叉树最近公共祖先}

\textbf{题目说明。} LeetCode 236 二叉树最近公共祖先 的题面入口要先还原成具体任务：普通树没有有序性，只能依靠祖先关系。

\textbf{样例入口。} 拿普通二叉树根 3，p=5、q=1 手算，左右子树分别找到一个目标，根 3 就是最近公共祖先；若两个目标都在左子树，答案应由左子树返回。

\textbf{暴力基线。} 慢版本在每个节点重新扫描子树或路径，先求出局部答案。重复扫描暴露了真正状态：每个节点只需要从孩子或父亲那里拿到少量信息。放到本题就是：先按“普通树没有有序性，只能依靠祖先关系”完整试慢版本；样例里已经能看到“规律是后序返回是否找到目标或祖先，不能用 BST 的大小关系”，所以优化前必须先能说清慢版本枚举了哪些候选。

\textbf{暴力过程。} 拿普通二叉树根 3，p=5、q=1 手算，左右子树分别找到一个目标，根 3 就是最近公共祖先；若两个目标都在左子树，答案应由左子树返回。

\textbf{重复工作。} 题面模型：普通树没有有序性，只能依靠祖先关系。暴力版的慢点要落到本题：慢版本会在树上重复确认“普通树没有有序性，只能依靠祖先关系”。样例规律是“规律是后序返回是否找到目标或祖先，不能用 BST 的大小关系”，说明父子之间真正需要传递的是高度、上下界、命中状态、层号或两类选择值，而不是反复扫描整棵子树。后面的优化只消掉这类重复工作，不能改变答案集合。

\textbf{优化条件。} 本题的关键观察是：父指针集合或后序返回目标命中状态都可行。这句话必须能翻译成可维护状态；如果题目条件改变后观察不再成立，就不能继续套原来的写法。

\textbf{相邻状态。} 规律是后序返回是否找到目标或祖先，不能用 BST 的大小关系。把这个特例继续拆开看：节点向父亲返回的量和全局答案通常不是同一个量。先写叶子或空节点的返回值，再写父节点如何合并左右孩子，递归式才有边界基础。

\textbf{状态复用。} 把对子树的重复扫描压成返回值或队列状态；每个节点只合并孩子给出的稳定信息。本题落点是：规律是后序返回是否找到目标或祖先，不能用 BST 的大小关系。手算时只改被新输入影响的那一小部分，而不是回头重算完整候选。

\textbf{指针和字段。} 状态字段要能说清：节点指针、传入约束或返回值、当前答案、层号或路径状态；空节点的返回值必须和状态定义匹配。手算时按这个协议跟踪：拿普通二叉树根 3，p=5、q=1 手算，左右子树分别找到一个目标，根 3 就是最近公共祖先；若两个目标都在左子树，答案应由左子树返回。然后按协议复查：手算时从叶子开始写返回值，或者从根开始写传入约束。不要只写访问顺序，要写每条边上传递的信息。

\textbf{代码落点。} 递归版本先处理空节点；迭代版本的栈元素要带完整状态；全负值、重复值和极端深树要单独测。关键代码行必须能逐句翻译成“旧状态怎样变成新状态”，否则就只是模板。

\textbf{正确性。} 证明从子树归纳或层序不变量开始：孩子状态正确时父节点合并正确；若用队列，当前轮队列必须正好保存当前层节点。

\textbf{边界实验。} 先用目标不存在、p 等于 q、一个是另一个祖先、比较节点地址做反例检查。实验时覆盖目标不存在、p 等于 q、一个是另一个祖先、比较节点地址，尤其关注空节点、单节点、链式树和全负值。递归版本和显式栈版本可以互相对拍。

\textbf{复杂度变化。} 暴力在每个节点重复扫描子树或反复寻找层边界可能退化到平方级；后序汇总、显式栈或层序遍历让每个节点处理常数次，通常是线性时间。

\textbf{迁移追问。} 如果题目改成“多次 LCA 查询可预处理倍增祖先表”，先判断原状态是否仍然覆盖新答案集合，再决定是微调边界、改 value 语义，还是换算法。

\subsection{LeetCode 1044 最长重复子串}

\textbf{题目说明。} LeetCode 1044 最长重复子串 的题面入口要先还原成具体任务：重复子串长度具有可行性单调：存在长度 L 的重复，则更短长度也存在。

\textbf{样例入口。} 拿 banana 手算，重复子串 ana 出现两次，长度 3 可行；若长度 4 不可行，则更长也不可能。

\textbf{暴力基线。} 慢版本按顺序扫描下标、逐个试答案值，或直接检查候选直到遇到目标边界。它先保证候选空间覆盖完整，但每次只排除一个候选，浪费了题目给出的顺序、比较或可行性边界。放到本题就是：先按“重复子串长度具有可行性单调：存在长度 L 的重复，则更短长度也存在”完整试慢版本；样例里已经能看到“规律是二分重复子串长度，用滚动哈希检查同长度子串是否重复，哈希相等时还要防冲突”，所以优化前必须先能说清慢版本枚举了哪些候选。

\textbf{暴力过程。} 拿 banana 手算，重复子串 ana 出现两次，长度 3 可行；若长度 4 不可行，则更长也不可能。

\textbf{重复工作。} 题面模型：重复子串长度具有可行性单调：存在长度 L 的重复，则更短长度也存在。暴力版的慢点要落到本题：慢版本会按顺序试“重复子串长度具有可行性单调：存在长度 L 的重复，则更短长度也存在”里的候选；而样例规律已经指出“规律是二分重复子串长度，用滚动哈希检查同长度子串是否重复，哈希相等时还要防冲突”。一次比较或一次可行性检查能丢掉一整段候选，继续线性试就是浪费。后面的优化只消掉这类重复工作，不能改变答案集合。

\textbf{优化条件。} 本题的关键观察是：二分长度，用滚动哈希检查是否存在两个相同长度子串哈希相同，并在冲突时确认字符串。这句话必须能翻译成可维护状态；如果题目条件改变后观察不再成立，就不能继续套原来的写法。

\textbf{相邻状态。} 规律是二分重复子串长度，用滚动哈希检查同长度子串是否重复，哈希相等时还要防冲突。把这个特例继续拆开看：每次比较 mid 之后，必须能指出哪一侧整体不可能包含目标，或哪一侧整体已经满足可行性。二分的核心不是 mid 公式，而是被丢弃区间确实不含答案。

\textbf{状态复用。} 把逐个试候选改成维护答案所在区间；边界谓词、可行性检查或局部比较给出分支方向，每个分支都必须能排除一侧候选。本题落点是：规律是二分重复子串长度，用滚动哈希检查同长度子串是否重复，哈希相等时还要防冲突。手算时只改被新输入影响的那一小部分，而不是回头重算完整候选。

\textbf{指针和字段。} 状态字段要能说清：left、right、mid、mid 处比较值或检查返回值、答案区间含义、返回值语义；每一轮更新都要保留目标位置或第一个可行值。手算时按这个协议跟踪：拿 banana 手算，重复子串 ana 出现两次，长度 3 可行；若长度 4 不可行，则更长也不可能。然后按协议复查：手算时列出 left、mid、right，以及 mid 处比较结果、可行性检查结果或局部趋势。每一步都要写出被丢弃的一侧为什么不含答案。

\textbf{代码落点。} 检查函数或比较分支单独写清，mid 不溢出，循环退出条件和返回值配套；每个分支都要解释丢弃区间。关键代码行必须能逐句翻译成“旧状态怎样变成新状态”，否则就只是模板。

\textbf{正确性。} 证明从排除一侧开始：答案空间题证明谓词单调，局部比较题证明保留下来的区间仍含目标；不能只靠经验猜测左右边界。

\textbf{边界实验。} 先用全相同字符、无重复、哈希冲突、长度边界、模数选择做反例检查。实验时覆盖全相同字符、无重复、哈希冲突、长度边界、模数选择，逐轮打印 left、mid、right、比较值或检查结果。只要有一次被排除区间仍含答案，二分不变量就错了。

\textbf{复杂度变化。} 暴力逐个试候选是线性或和值域成正比；二分每次排除一半候选，复杂度变成对数级，答案二分还要乘上单次检查成本。

\textbf{迁移追问。} 如果题目改成“也可用后缀数组或后缀自动机，适合系统学习字符串算法”，先判断原状态是否仍然覆盖新答案集合，再决定是微调边界、改 value 语义，还是换算法。

\subsection{LeetCode 337 打家劫舍 III}

\textbf{题目说明。} LeetCode 337 打家劫舍 III 的题面入口要先还原成具体任务：树上相邻节点不能同时选，每个节点需要偷和不偷两个状态。

\textbf{样例入口。} 拿树 [3, 2, 3, null, 3, null, 1] 手算，偷根 3 时不能偷两个孩子，只能接孙子；不偷根时可以分别取左右孩子的最优。

\textbf{暴力基线。} 慢版本在每个节点重新扫描子树或路径，先求出局部答案。重复扫描暴露了真正状态：每个节点只需要从孩子或父亲那里拿到少量信息。放到本题就是：先按“树上相邻节点不能同时选，每个节点需要偷和不偷两个状态”完整试慢版本；样例里已经能看到“规律是每个节点返回两个状态：偷当前和不偷当前，父节点再组合”，所以优化前必须先能说清慢版本枚举了哪些候选。

\textbf{暴力过程。} 拿树 [3, 2, 3, null, 3, null, 1] 手算，偷根 3 时不能偷两个孩子，只能接孙子；不偷根时可以分别取左右孩子的最优。

\textbf{重复工作。} 题面模型：树上相邻节点不能同时选，每个节点需要偷和不偷两个状态。暴力版的慢点要落到本题：慢版本会在树上重复确认“树上相邻节点不能同时选，每个节点需要偷和不偷两个状态”。样例规律是“规律是每个节点返回两个状态：偷当前和不偷当前，父节点再组合”，说明父子之间真正需要传递的是高度、上下界、命中状态、层号或两类选择值，而不是反复扫描整棵子树。后面的优化只消掉这类重复工作，不能改变答案集合。

\textbf{优化条件。} 本题的关键观察是：后序汇总：偷当前则孩子不能偷，不偷当前则孩子可偷可不偷。这句话必须能翻译成可维护状态；如果题目条件改变后观察不再成立，就不能继续套原来的写法。

\textbf{相邻状态。} 规律是每个节点返回两个状态：偷当前和不偷当前，父节点再组合。把这个特例继续拆开看：节点向父亲返回的量和全局答案通常不是同一个量。先写叶子或空节点的返回值，再写父节点如何合并左右孩子，递归式才有边界基础。

\textbf{状态复用。} 把对子树的重复扫描压成返回值或队列状态；每个节点只合并孩子给出的稳定信息。本题落点是：规律是每个节点返回两个状态：偷当前和不偷当前，父节点再组合。手算时只改被新输入影响的那一小部分，而不是回头重算完整候选。

\textbf{指针和字段。} 状态字段要能说清：节点指针、传入约束或返回值、当前答案、层号或路径状态；空节点的返回值必须和状态定义匹配。手算时按这个协议跟踪：拿树 [3, 2, 3, null, 3, null, 1] 手算，偷根 3 时不能偷两个孩子，只能接孙子；不偷根时可以分别取左右孩子的最优。然后按协议复查：手算时从叶子开始写返回值，或者从根开始写传入约束。不要只写访问顺序，要写每条边上传递的信息。

\textbf{代码落点。} 递归版本先处理空节点；迭代版本的栈元素要带完整状态；全负值、重复值和极端深树要单独测。关键代码行必须能逐句翻译成“旧状态怎样变成新状态”，否则就只是模板。

\textbf{正确性。} 证明从子树归纳或层序不变量开始：孩子状态正确时父节点合并正确；若用队列，当前轮队列必须正好保存当前层节点。

\textbf{边界实验。} 先用空树、负值是否可能、链式树、递归深度做反例检查。实验时覆盖空树、负值是否可能、链式树、递归深度，尤其关注空节点、单节点、链式树和全负值。递归版本和显式栈版本可以互相对拍。

\textbf{复杂度变化。} 暴力在每个节点重复扫描子树或反复寻找层边界可能退化到平方级；后序汇总、显式栈或层序遍历让每个节点处理常数次，通常是线性时间。

\textbf{迁移追问。} 如果题目改成“可用显式栈后序保存节点状态，避免调用栈”，先判断原状态是否仍然覆盖新答案集合，再决定是微调边界、改 value 语义，还是换算法。

\subsection{LeetCode 486 预测赢家}

\textbf{题目说明。} LeetCode 486 预测赢家 的题面入口要先还原成具体任务：两人从区间两端取数，当前选择会影响对手面对的剩余区间。

\textbf{样例入口。} 拿 nums=[1, 5, 2] 手算，先手若拿 1，后手会拿 5；若先手拿 2，后手仍拿 5，先手最终都会输。dp[left][right] 保存当前玩家比对手最多多拿多少分。

\textbf{暴力基线。} 慢版本先写递归搜索树：节点表示处理位置、剩余资源和当前选择。只要不同路径反复到达同一个节点，就说明需要把这个节点命名成 DP 状态。放到本题就是：先按“两人从区间两端取数，当前选择会影响对手面对的剩余区间”完整试慢版本；样例里已经能看到“规律是选左得到 nums[left]-dp[left+1][right]，选右得到 nums[right]-dp[left][right-1]”，所以优化前必须先能说清慢版本枚举了哪些候选。

\textbf{暴力过程。} 拿 nums=[1, 5, 2] 手算，先手若拿 1，后手会拿 5；若先手拿 2，后手仍拿 5，先手最终都会输。dp[left][right] 保存当前玩家比对手最多多拿多少分。

\textbf{重复工作。} 题面模型：两人从区间两端取数，当前选择会影响对手面对的剩余区间。暴力版的慢点要落到本题：慢版本会在“两人从区间两端取数，当前选择会影响对手面对的剩余区间”的选择树里反复到达相同参数。样例规律是“规律是选左得到 nums[left]-dp[left+1][right]，选右得到 nums[right]-dp[left][right-1]”，说明这个参数组合可以命名成状态，算过之后后面直接复用。后面的优化只消掉这类重复工作，不能改变答案集合。

\textbf{优化条件。} 本题的关键观察是：dp[left][right] 表示当前玩家相对对手最多能赢多少分，取左或取右后减去子区间优势。这句话必须能翻译成可维护状态；如果题目条件改变后观察不再成立，就不能继续套原来的写法。

\textbf{相邻状态。} 规律是选左得到 nums[left]-dp[left+1][right]，选右得到 nums[right]-dp[left][right-1]。把这个特例继续拆开看：先问暴力搜索里哪些不同路径到达同一个后续问题，再给这个后续问题命名为状态。状态必须包含未来决策需要的全部信息，转移只从更小且已经算清的状态来。

\textbf{状态复用。} 把重复递归节点命名为状态；遍历顺序只是在保证依赖状态已经算完。本题落点是：规律是选左得到 nums[left]-dp[left+1][right]，选右得到 nums[right]-dp[left][right-1]。手算时只改被新输入影响的那一小部分，而不是回头重算完整候选。

\textbf{指针和字段。} 状态字段要能说清：状态下标、状态值语义、初始化、转移来源、遍历顺序；空间压缩时还要指出读到的是旧值还是新值。手算时按这个协议跟踪：拿 nums=[1, 5, 2] 手算，先手若拿 1，后手会拿 5；若先手拿 2，后手仍拿 5，先手最终都会输。dp[left][right] 保存当前玩家比对手最多多拿多少分。然后按协议复查：手算时先写一个很小的 DP 表或记忆化搜索记录。每个格子旁边写它来自哪些更小状态。

\textbf{代码落点。} 先写未压缩版本校验转移；压缩前后用小表对拍；不可达哨兵参与加法前要检查溢出。关键代码行必须能逐句翻译成“旧状态怎样变成新状态”，否则就只是模板。

\textbf{正确性。} 证明从最后一步开始：任意最优解的最后一步都在转移枚举中，转移来源已经由更小状态正确计算。

\textbf{边界实验。} 先用单元素、两个元素、总分相等、负数若存在、区间遍历顺序做反例检查。实验时覆盖单元素、两个元素、总分相等、负数若存在、区间遍历顺序，先打印完整 DP 表，再比较空间压缩版本。若压缩版不同，优先检查遍历顺序和旧值保存。

\textbf{复杂度变化。} 暴力递归会重复计算相同子问题，常见指数级；DP 把每个状态算一次，时间是状态数乘单个转移成本，空间是状态表大小或压缩后的大小。

\textbf{迁移追问。} 如果题目改成“若要输出具体策略，需要记录每个区间选择的端点”，先判断原状态是否仍然覆盖新答案集合，再决定是微调边界、改 value 语义，还是换算法。

\subsection{LeetCode 200 岛屿数量}

\textbf{题目说明。} LeetCode 200 岛屿数量 的题面入口要先还原成具体任务：陆地格子是节点，四方向相邻是边，岛屿是连通块。

\textbf{样例入口。} 拿一个 2x2 小岛手算，从一个陆地出发能沿上下左右标记同一块岛。若不标记访问，会重复计数；若对角线也连通，会误合并。

\textbf{暴力基线。} 暴力可以对每个陆地格子单独搜索它所属的连通块，再用集合去重；这样同一块岛会被重复搜索很多次。

\textbf{暴力过程。} 以 2x2 全是陆地为例，从左上搜索会访问四个格子；如果之后从右上又重新搜索，就把同一座岛重复数了一遍。

\textbf{重复工作。} 题面模型：陆地格子是节点，四方向相邻是边，岛屿是连通块。暴力版的慢点要落到本题：浪费在已确认属于某个岛的陆地没有被标记。每块岛只需要在第一次遇到时完整淹没或标记。后面的优化只消掉这类重复工作，不能改变答案集合。

\textbf{优化条件。} 本题的关键观察是：扫描遇到未访问陆地就启动搜索并标记整块。这句话必须能翻译成可维护状态；如果题目条件改变后观察不再成立，就不能继续套原来的写法。

\textbf{相邻状态。} 规律是每发现一个未访问陆地，答案加一，并用 DFS/BFS 标记与它四连通的全部陆地。把这个特例继续拆开看：状态节点不一定等于原始位置，还可能包含钥匙、剩余资源、时间或访问集合。visited 只能合并未来能力完全相同的状态，否则会把合法路径剪掉。

\textbf{状态复用。} 扫描网格。遇到未访问陆地，答案加一，然后 BFS/DFS 把与它四方向连通的全部陆地标记为已访问。手算时只改被新输入影响的那一小部分，而不是回头重算完整候选。

\textbf{指针和字段。} 状态字段要能说清：row/col 是当前格子，visited 或原地改写表示已经归属某座岛，directions 是四方向邻居。手算时按这个协议跟踪：从左上陆地入队，依次扩展右边和下边陆地，最终整块连通区域都标记；扫描到这些格子时直接跳过。

\textbf{代码落点。} 关键是只看上下左右，不看对角线。边界判断集中写，访问标记在入队时设置，避免重复入队。关键代码行必须能逐句翻译成“旧状态怎样变成新状态”，否则就只是模板。

\textbf{正确性。} 一次搜索会访问且只访问与起点四连通的全部陆地；外层每次启动搜索都对应一块此前未计数的新岛。

\textbf{边界实验。} 先用空网格、全水、全陆、斜对角不连通、是否允许修改输入做反例检查。实验时覆盖空网格、全水、全陆、斜对角不连通、是否允许修改输入，并打印入队时的完整状态。若同一位置但资源不同的状态被合并，很多图题会出现隐蔽错误。

\textbf{复杂度变化。} 重复搜索最坏 O(\ensuremath{(mn)^{2}})，标记搜索 O(mn) 时间，空间 O(mn) 或递归/队列大小。

\textbf{迁移追问。} 如果题目改成“也可用并查集合并相邻陆地，适合动态岛屿扩展”，先判断原状态是否仍然覆盖新答案集合，再决定是微调边界、改 value 语义，还是换算法。

\subsection{LeetCode 312 戳气球}

\textbf{题目说明。} LeetCode 312 戳气球 的题面入口要先还原成具体任务：先戳哪个会让邻居持续变化，改成最后戳某个气球后区间边界才稳定。

\textbf{样例入口。} 拿 [3, 1, 5] 手算，若先戳 1，左右邻居会变；若改问某个区间里最后戳谁，最后一刻左右边界固定，收益可拆成左右两个子区间。

\textbf{暴力基线。} 慢版本先写递归搜索树：节点表示处理位置、剩余资源和当前选择。只要不同路径反复到达同一个节点，就说明需要把这个节点命名成 DP 状态。放到本题就是：先按“先戳哪个会让邻居持续变化，改成最后戳某个气球后区间边界才稳定”完整试慢版本；样例里已经能看到“规律是区间 DP 枚举最后戳的 mid，而不是第一个戳的气球”，所以优化前必须先能说清慢版本枚举了哪些候选。

\textbf{暴力过程。} 拿 [3, 1, 5] 手算，若先戳 1，左右邻居会变；若改问某个区间里最后戳谁，最后一刻左右边界固定，收益可拆成左右两个子区间。

\textbf{重复工作。} 题面模型：先戳哪个会让邻居持续变化，改成最后戳某个气球后区间边界才稳定。暴力版的慢点要落到本题：慢版本会在“先戳哪个会让邻居持续变化，改成最后戳某个气球后区间边界才稳定”的选择树里反复到达相同参数。样例规律是“规律是区间 DP 枚举最后戳的 mid，而不是第一个戳的气球”，说明这个参数组合可以命名成状态，算过之后后面直接复用。后面的优化只消掉这类重复工作，不能改变答案集合。

\textbf{优化条件。} 本题的关键观察是：区间 DP 定义开区间内全部戳完的最大收益，枚举最后一个被戳的气球作为切分点。这句话必须能翻译成可维护状态；如果题目条件改变后观察不再成立，就不能继续套原来的写法。

\textbf{相邻状态。} 规律是区间 DP 枚举最后戳的 mid，而不是第一个戳的气球。把这个特例继续拆开看：先问暴力搜索里哪些不同路径到达同一个后续问题，再给这个后续问题命名为状态。状态必须包含未来决策需要的全部信息，转移只从更小且已经算清的状态来。

\textbf{状态复用。} 把重复递归节点命名为状态；遍历顺序只是在保证依赖状态已经算完。本题落点是：规律是区间 DP 枚举最后戳的 mid，而不是第一个戳的气球。手算时只改被新输入影响的那一小部分，而不是回头重算完整候选。

\textbf{指针和字段。} 状态字段要能说清：状态下标、状态值语义、初始化、转移来源、遍历顺序；空间压缩时还要指出读到的是旧值还是新值。手算时按这个协议跟踪：拿 [3, 1, 5] 手算，若先戳 1，左右邻居会变；若改问某个区间里最后戳谁，最后一刻左右边界固定，收益可拆成左右两个子区间。然后按协议复查：手算时先写一个很小的 DP 表或记忆化搜索记录。每个格子旁边写它来自哪些更小状态。

\textbf{代码落点。} 先写未压缩版本校验转移；压缩前后用小表对拍；不可达哨兵参与加法前要检查溢出。关键代码行必须能逐句翻译成“旧状态怎样变成新状态”，否则就只是模板。

\textbf{正确性。} 证明从最后一步开始：任意最优解的最后一步都在转移枚举中，转移来源已经由更小状态正确计算。

\textbf{边界实验。} 先用补一的虚拟边界、区间长度遍历顺序、空区间收益为零、乘积可能较大做反例检查。实验时覆盖补一的虚拟边界、区间长度遍历顺序、空区间收益为零、乘积可能较大，先打印完整 DP 表，再比较空间压缩版本。若压缩版不同，优先检查遍历顺序和旧值保存。

\textbf{复杂度变化。} 暴力递归会重复计算相同子问题，常见指数级；DP 把每个状态算一次，时间是状态数乘单个转移成本，空间是状态表大小或压缩后的大小。

\textbf{迁移追问。} 如果题目改成“很多区间 DP 都要换成最后一步思考，避免中间操作改变边界”，先判断原状态是否仍然覆盖新答案集合，再决定是微调边界、改 value 语义，还是换算法。

\begin{principlebox}{本节复盘方式}
不要只看最终算法名。每道题复盘时先遮住优化版，口述暴力枚举对象；再指出相邻窗口、历史前缀、候选区间、搜索状态或子问题里哪一部分被重复处理；最后用一个边界样例验证状态更新顺序。能讲完这三步，才算真正掌握本章案例。
\end{principlebox}
